% 分布式数据库系统 往年试题参考解析
% 用 xelatex 编译两遍（第二遍生成页码引用）。需要字体：Noto Serif CJK SC、Noto Sans CJK SC、DejaVu Sans Mono、Latin Modern Roman。
\documentclass[10pt,a4paper]{article}
\usepackage[a4paper,top=1.7cm,bottom=1.6cm,left=1.6cm,right=1.6cm,headsep=0.4cm,footskip=0.7cm]{geometry}
\usepackage{fontspec}
\setmainfont{Noto Serif CJK SC}[BoldFont={Noto Serif CJK SC Bold},ItalicFont={Noto Serif CJK SC},ItalicFeatures={FakeSlant=0.18}]
\setsansfont{Noto Sans CJK SC}[BoldFont={Noto Sans CJK SC Bold}]
\setmonofont{DejaVu Sans Mono}[Scale=0.86]
\newfontfamily\zhmono{Noto Sans CJK SC}[Scale=0.86]
\XeTeXlinebreaklocale "zh"
\XeTeXlinebreakskip = 0pt plus 1pt minus 0.1pt
\linespread{1.18}
\setlength{\parindent}{0pt}
\setlength{\parskip}{3pt plus 1pt}

\usepackage{amsmath,amssymb,mathtools}
\usepackage{xcolor}
\usepackage{pifont}
\usepackage{array,booktabs,tabularx,multirow,makecell}
\usepackage{enumitem}
\setlist{nosep,leftmargin=1.6em,itemsep=1pt}
\usepackage{fancyvrb}
\usepackage{needspace}
\usepackage{tikz}
\usetikzlibrary{arrows.meta,positioning,calc,shapes.misc,decorations.pathreplacing,fit}
\usepackage{forest}
\usepackage[most]{tcolorbox}
\usepackage{fancyhdr}
\usepackage{titlesec}
\usepackage{hyperref}
\hypersetup{hidelinks,bookmarksnumbered}

% 标题
\titleformat{\section}{\sffamily\bfseries\Large}{%
  \colorbox{black}{\color{white}\,第\thesection 部分\,}}{0.6em}{}[\vspace{-2pt}\rule{\linewidth}{1.2pt}]
\titlespacing*{\section}{0pt}{0pt}{6pt}
\titleformat{\subsection}{\sffamily\bfseries\large}{\thesubsection}{0.5em}{}
\titlespacing*{\subsection}{0pt}{8pt}{3pt}

% 页眉页脚
\pagestyle{fancy}
\fancyhf{}
\renewcommand{\headrulewidth}{0.5pt}
\fancyhead[L]{\sffamily\small 分布式数据库系统}
\fancyhead[R]{\sffamily\small\nouppercase{\leftmark}}
\fancyfoot[C]{\sffamily\small — \thepage\ —}
\renewcommand{\sectionmark}[1]{\markboth{第\thesection 部分\ #1}{}}

% 记号
\newcommand{\sj}{\ltimes}
\newcommand{\jn}{\bowtie}
\newcommand{\jnon}[1]{\mathbin{\bowtie_{\text{#1}}}}
\newcommand{\sjon}[1]{\mathbin{\ltimes_{\text{#1}}}}
\newcommand{\R}[1]{\mathrm{#1}}
\newcommand{\sel}[1]{\sigma_{\text{#1}}}
\newcommand{\proj}[1]{\Pi_{\text{#1}}}
\newcommand{\cir}[1]{\ding{\numexpr171+#1\relax}}
\newcommand{\bcir}[1]{\ding{\numexpr181+#1\relax}}
\newcommand{\kj}[1]{\textsf{\small[课件#1]}}
\newcommand{\zh}[1]{{\zhmono #1}}
\newcommand{\cross}{\ding{55}}
\newfontfamily\lmr{Latin Modern Roman}
\newcommand{\q}[1]{{\lmr'}#1{\lmr'}}
\newcommand{\mand}{\ensuremath{\,\wedge\,}}

% 框（黑白打印）
\tcbset{
  base/.style={enhanced,sharp corners,boxsep=1.5pt,left=5pt,right=5pt,top=3pt,bottom=3pt,
    before skip=6pt,after skip=6pt,fonttitle=\sffamily\bfseries\small},
  tab/.style={attach boxed title to top left={xshift=4mm,yshift*=-\tcboxedtitleheight/2},top=3.4mm},
}
% 做法：左侧粗黑条
\newtcolorbox{flow}[1][]{base,tab,colback=white,colframe=black,boxrule=0.6pt,leftrule=3.5pt,
  title={#1},coltitle=white,colbacktitle=black,boxed title style={sharp corners,boxrule=0pt}}
% 参考答卷：双线，不跨页
\newtcolorbox{ans}[1][]{base,tab,colback=white,colframe=black,boxrule=0.5pt,
  borderline={0.5pt}{1.6pt}{black},
  title={参考答卷\ifx\relax#1\relax\else：#1\fi},coltitle=black,colbacktitle=white,
  boxed title style={sharp corners,boxrule=0.5pt}}
% 注意：虚线框
\newtcolorbox{pit}[1][]{base,tab,colback=white,frame hidden,
  borderline={0.7pt}{0pt}{black,dashed},
  title={注意\ifx\relax#1\relax\else：#1\fi},coltitle=black,colbacktitle=white,
  boxed title style={sharp corners,boxrule=0pt}}
% 其他：细框
\newtcolorbox{var}[1][]{base,tab,colback=white,colframe=black!60,boxrule=0.4pt,
  title={#1},coltitle=black,colbacktitle=white,
  boxed title style={sharp corners,boxrule=0.4pt,colframe=black!60}}
% 原题
\newtcolorbox{prob}[1][]{base,colback=white,colframe=black,boxrule=0.3pt,
  title={#1},coltitle=white,colbacktitle=black!75,fontupper=\small}

% 问答（同一组问答不拆到两页）
\newcommand{\Qlab}{\colorbox{black}{\color{white}\sffamily\bfseries\footnotesize 问}}
\newcommand{\Alab}{\fcolorbox{black}{white}{\sffamily\bfseries\footnotesize 答}}
\newcommand{\Q}[1]{\par\Needspace{6\baselineskip}\hangindent=2.1em\hangafter=1\noindent\makebox[2.1em][l]{\Qlab}{\sffamily #1}\par\nopagebreak}
\newcommand{\A}[1]{\par\hangindent=2.1em\hangafter=1\noindent\makebox[2.1em][l]{\Alab}#1\par\vspace{1pt}}
\newenvironment{dlg}{\par\vspace{2pt}\begingroup\small}{\endgroup\par\vspace{2pt}}

% 树
\forestset{
  qt/.style={for tree={font=\footnotesize,inner sep=0.6pt,outer sep=0pt,l=0pt,l sep=7pt,s sep=5pt,
    parent anchor=south,child anchor=north,edge={thin},align=center,base=bottom}},
  ft/.style={for tree={font=\footnotesize,draw,rounded corners=1pt,inner sep=1.5pt,l=0pt,l sep=15pt,s sep=6pt,
    parent anchor=south,child anchor=north,edge={-{Stealth[length=3pt]}}}},
  h/.style={edge label={node[midway,fill=white,inner sep=0.5pt,font=\scriptsize\itshape]{h}}},
  v/.style={edge label={node[midway,fill=white,inner sep=0.5pt,font=\scriptsize\itshape]{v}}},
  cut/.style={edge={dashed,black!45},for tree={font=\small\color{black!45},edge={dashed,black!45}}},
}

% 代码
\DefineVerbatimEnvironment{code}{Verbatim}{commandchars=\\\{\},fontsize=\small,frame=leftline,framerule=1.2pt,xleftmargin=4pt,baselinestretch=1.05}

% 题目与答卷放在同一页
\newenvironment{keep}{\par\noindent\begin{minipage}{\linewidth}\setlength{\parskip}{3pt plus 1pt}}{\end{minipage}\par\vspace{2pt}}

\begin{document}
\thispagestyle{empty}
\begin{center}
{\sffamily\bfseries\LARGE 分布式数据库系统\ \ 往年试题参考解析}\par\vspace{3pt}
{\small 2018、2019、2020、2023 年试题，以及 2004–2011 年的概念题}
\end{center}

{\small 题目按类型分成七部分。每类题先写做法，然后用问答讲清每一步的理由，灰底双线框是参考答卷，虚线框是容易出错的地方。“课件 3-21”指打印版 PPT 每页右上角的编号。}

\subsection*{按题型查找}
\vspace{-2pt}
{\small
\renewcommand{\arraystretch}{1.05}
\begin{tabularx}{\linewidth}{@{}>{\raggedright\arraybackslash}X >{\raggedright\arraybackslash}p{3.7cm} l l@{}}
\toprule
\textbf{题目里出现} & \textbf{题型} & \textbf{本文} & \textbf{课件} \\
\midrule
“按需求分片”“分片定义”“分片树”“分片类型” & 水平 / 垂直 / 导出 / 混合 & \S1.1\,p.\pageref{s:frag} & 3-21--55 \\
“分配设计”“分配类型” & 全分割 / 全复制 / 部分复制 & \S1.1\,p.\pageref{s:frag} & 3-57--63 \\
“写出关系代数和 SQL”“全局查询树” & 查询 → 关系代数 → 查询树 & \S1.2\,p.\pageref{s:opt} & 5-21--23, 48 \\
“全局优化，写出中间过程” & 选择、投影下推 & \S1.2\,p.\pageref{s:opt} & 5-49--55 \\
“分片优化”“片段查询树” & 本地化与剪枝 & \S1.2\,p.\pageref{s:opt} & 5-57--69 \\
“内存存放 5 块”“每秒读写 25 块” & I/O 代价 & \S2.1\,p.\pageref{s:io} & 5-11--18 \\
“S1、S2、S3，选传输代价最小的执行场地” & 选择执行场地 & \S2.2\,p.\pageref{s:site} & 6-5--7 \\
“全连接”“优化的查询策略”“通信代价”，给了 Card、Length、Val & 直接连接与半连接 & \S2.3\,p.\pageref{s:semi} & 6-15--39 \\
“分解为子事务”“进程代理模型” & 子事务与伪代码 & \S3.1\,p.\pageref{s:proc} & 7-19--26 \\
“集中式两阶段提交”“协调者和参与者” & 2PC 流程与日志 & \S3.2\,p.\pageref{s:2pc} & 7-29--34, 47 \\
“某场地出故障”“丢失 commit 命令”“转账没有成功” & 2PC 的故障处理 & \S3.3\,p.\pageref{s:fail} & 8-42--55 \\
“冲突对”“局部可串行”“全局可串行”“构建历程” & 可串行化 & \S4.1\,p.\pageref{s:cc} & 9-10--18 \\
“两阶段封锁协议” & 2PL & \S4.5\,p.\pageref{s:2pl} & 9-24--33 \\
日志 $B_1\,D^1_1\,C_1\,K\dots$，“undo”“redo”“检查点” & 故障恢复 & \S5\,p.\pageref{s:rec} & 8-13--40 \\
“主副本”“同步复制协议”“如何保证一致性” & 主从同步复制 & \S6.1\,p.\pageref{s:rep} & 10-11--28 \\
“Paxos、反熵、NWR 的异同” & 一致性协议比较 & \S6.2\,p.\pageref{s:paxos} & 10-38--54 \\
ACID、集中式与分布式的区别、区块链、学习收获 & 概念题 & \S7\,p.\pageref{s:concept} & 7-9--14, 1-52 \\
\bottomrule
\end{tabularx}}

\subsection*{记号}
\vspace{-2pt}
{\small
\begin{tabularx}{\linewidth}{@{}l X l X@{}}
\toprule
$\sel{F}(R)$ & 选择，留下满足 $F$ 的元组 & $R\jn S$ & 连接，课件写作 $R\infty S$ \\
$\proj{A}(R)$ & 投影，留下属性 $A$ & $R\sj S$ & 半连接，课件写作 $R\propto S$，等于 $\proj{Attr(R)}(R\jn S)$ \\
h、v & 分片树的边：水平、垂直 & $\cup$ & 并。水平片段用 $\cup$ 重构，垂直片段用 $\jn$ 重构\\
Card$(R)$ & 元组数 & Length$(R)$ & 一个元组的字节数 \\
Size$(R)$ & $\text{Card}(R)\times\text{Length}(R)$ & Val$(R,A)$ & 属性 $A$ 不同取值的个数 \\
$\rho$ & 选择度，满足条件的元组所占比例 & $C_0,C_1$ & 一次传输的启动代价，单位数据传输代价 \\
$B_i, C_i, A_i$ & 事务 $T_i$ 开始、提交、废弃 & $D^j_i$ & $T_i$ 的第 $j$ 个数据操作；$K$ 是检查点 \\
$R_i(x), W_i(x)$ & $T_i$ 读、写数据项 $x$ & $T_i<T_j$ & 在等价的串行顺序中 $T_i$ 在 $T_j$ 前 \\
\bottomrule
\end{tabularx}}

\vspace{4pt}
{\small 计算题和设计题的答题顺序一般是：先写假设（数据均匀分布、单位、属性取值的含义），再写依据的定义或公式，然后分步计算或作图，每步注明用了哪条规则，最后写结论。题目要求“写出中间过程”时，每次等价变换都要单独写一行或画一棵树。}

\clearpage
\section{分片设计、分配设计与查询优化}
{\small 每年第 1 题都是这一类，约 20 分。常见的五问是：分片定义、分片树和分片类型；分配设计和分配类型；SQL、关系代数和全局查询树；全局优化；分片优化。}

\subsection{分片设计与分配设计}\label{s:frag}
\begin{flow}[分片与分配设计的步骤]
\begin{enumerate}
\item 把每条需求写成“哪个关系的哪些数据存在哪个场地”。
\item 确定分片类型。按元组划分是水平分片，例如“各单位的员工”“涉密项目”；按属性划分是垂直分片，例如“属性 \{…\} 存在某场地”。划分元组的条件属性在另一个关系里时（“存在负责人所在场地”“存在发生借贷的分行”），用导出水平分片，写成半连接 $\sj$。先垂直再水平是混合分片。
\item 需求里有“其中满足 $C$ 的在某场地保存副本”时，把 $C$ 与 $\neg C$ 也作为分片条件再切一次。
\item 写分片定义。每个片段写成 $\sel{条件}(R)$、$\proj{属性}(R)$ 或 $R\sj S'$；同一层的条件互斥，并且覆盖全部取值。
\item 用一句话说明满足三条原则：完备性（片段的并等于原关系），可重构性（用 $\cup$ 或 $\jn$ 能还原），不相交性（水平片段两两相交为空，垂直片段之间只共享主键）。
\item 画分片树。根是全局关系，叶子是最终片段，边上标 h 或 v \kj{3-54}。
\item 分配设计：列出每个场地存放的片段。每个片段只存一处是全分割，每个片段在每个场地都有是全复制，其余是部分复制 \kj{3-57}。
\end{enumerate}
\end{flow}

\begin{dlg}
\Q{分片和分配有什么区别？}
\A{分片是逻辑划分，把全局关系分成片段；分配是物理放置，决定片段存在哪个场地 \kj{3-13}。片段是分配的最小单位，要放就整个放到某个场地，不能只放一部分。同一个片段存在多个场地，就是复制。}
\Q{“金额 $\ge$100 万的项目在总部保存副本”，为什么要因此多切一次？}
\A{副本只能以整个片段为单位。如果某个分支机构的项目都在片段 $P_1$ 里，而总部只要其中金额大的部分，就没法复制。把 $P_1$ 分成 $P_{11}$（$\ge$100）和 $P_{12}$（$<$100）以后，$P_{11}$ 放两个场地，$P_{12}$ 只放一个场地。题目里每出现一条副本要求，一般就要多一个分片条件。}
\Q{什么时候用导出水平分片？}
\A{划分依据的属性在别的关系里时。例如“企业信息存在与它有借贷关系的分行”：企业表里没有分行号，分行号 Bid 在贷款表里，所以用半连接把贷款表的划分带过来：$B_i=\R{Loanbusiness}\sjon{Lid}\sel{Bid=\q{Ai}}(\R{Loaninfo})$，即在 $A_i$ 有贷款的企业 \kj{3-36--38}。}
\Q{垂直分片为什么每个片段都要带主键？}
\A{垂直片段靠连接重构（$R=R_1\jn R_2$），连接需要公共属性，这个公共属性就是主键。课件的条件是 $\text{Attr}(R_i)\cap\text{Attr}(R_j)=\text{PK}(R)$ \kj{3-40}：主键以外的属性只能出现在一个片段里。}
\Q{分片条件怎么写？}
\A{同一层的条件两两矛盾，保证不相交；合起来恒为真，保证完备。成对写 $C$ 和 $\neg C$ 不会出错。按 Bid=A1/A2/A3 枚举时，要求 Bid 只有这几个取值，可以写一句“设 Bid 只取 A0–A3”。}
\end{dlg}

\begin{pit}
\begin{itemize}
\item 混合分片第二层的对象是垂直片段，要写 $\sel{Dno=0}(\R{Student_1})$，不写 $\sel{Dno=0}(\R{Student})$。
\item 分配类型按整个方案判断：有片段被复制，又不是所有片段都在所有场地，就是部分复制。
\item 不分片的关系也要写出来，例如“DEPT 不分片，存放在 S0”。
\end{itemize}
\end{pit}

\Needspace{14\baselineskip}
\subsection{查询优化：全局优化和分片优化}\label{s:opt}
\begin{flow}[查询优化的步骤]
\begin{enumerate}
\item SQL：FROM 写涉及的全局关系；WHERE 写连接条件和选择条件；SELECT 写要输出的属性。
\item 关系代数写成 $\proj{输出属性}\big(\sel{全部选择条件}(R\jn S\jn\cdots)\big)$。
\item 初始全局查询树：根是 $\Pi$，下面是 $\sigma$，再下面是 $\jn$，叶子是全局关系 \kj{5-48}。
\item 全局优化 \kj{5-49--55}，每一步注明依据：
 (a) 用选择串接律拆开 $\sel{$F_1\wedge F_2$}$；
 (b) 用选择与连接的分配律，把每个条件移到只含其属性的关系上方；
 (c) 用投影串接律和投影与连接的分配律，在每个叶子上方加 $\Pi$，只留输出属性和连接属性；
 (d) 画出优化后的全局查询树。
\item 分片优化 \kj{5-57--69}：
 (a) 数据本地化，把全局关系换成分片树，h 换成 $\cup$，v 换成 $\jn$；
 (b) 准则 1：把 $\sigma$ 移到每个片段上，与片段条件矛盾的置为空并删去；
 (c) 准则 3：把 $\jn$ 移到 $\cup$ 下面，$(A\cup B)\jn C=(A\jn C)\cup(B\jn C)$；准则 2：连接条件不可能满足的分支为空；
 (d) 准则 4：删去对查询没有贡献的垂直片段；
 (e) 画出最终的片段查询树，注明各片段所在场地。
\end{enumerate}
\end{flow}

{\small
\begin{tabularx}{\linewidth}{@{}l X@{}}
\toprule
\multicolumn{2}{@{}l}{\textbf{全局优化常用的等价变换规则} \kj{5-49--50}（写中间过程时在括号里注明）}\\
\midrule
选择串接律 & $\sel{$F_1$}(\sel{$F_2$}(R))=\sel{$F_1\wedge F_2$}(R)$ \\
选择与连接（笛卡尔积）的分配律 & 若 $F$ 只涉及 $R$ 的属性：$\sel{F}(R\jn S)=\sel{F}(R)\jn S$ \\
选择与并的分配律 & $\sel{F}(R\cup S)=\sel{F}(R)\cup\sel{F}(S)$，分片优化中把 $\sigma$ 移进 $\cup$ 用它\\
投影串接律 & $A\subseteq B$ 时 $\proj{A}(\proj{B}(R))=\proj{A}(R)$ \\
投影与连接的分配律 & $\proj{$A\cup B$}(R\jn S)=\proj{A}(R)\jn\proj{B}(S)$，$A$、$B$ 要包含连接属性 \\
\bottomrule
\end{tabularx}}

\begin{dlg}
\Q{为什么先把选择和投影往下移？}
\A{课件的原则是先做能缩小中间结果的运算 \kj{5-51}。$\sigma$、$\Pi$ 的代价是 $O(n)$，连接是 $O(n\log n)$，笛卡尔积是 $O(n^2)$。先用 $\sigma$ 减少元组、用 $\Pi$ 去掉用不到的属性，后面的连接就便宜得多。2023 年第三题里同一个查询的三种写法，最慢的约 22 小时，最快的 8 秒。}
\Q{投影下移时，为什么不能只保留 SELECT 里的属性？}
\A{上面还要做连接。把 $\proj{Sname}$ 移到 Student 上方时，如果连 Sno 也去掉，Student 就无法再和 SC 按 Sno 连接。每个叶子上方保留的是输出属性加上层连接要用的属性。}
\Q{条件 Dno=2 只写在 Student 上，DEPT 上也能加吗？}
\A{能。WHERE 里有 Student.Dno=DEPT.Dno，所以 Student.Dno=2 也意味着 DEPT.Dno=2。两边都先筛选，连接的两个输入都变小。2023 年的参考答案也是两边都加。}
\Q{全局优化和分片优化有什么区别？}
\A{全局优化只用代数变换规则，处理的是全局关系。分片优化把全局关系换成片段以后，可以根据片段的定义条件判断哪些片段查不出结果，直接删掉。最后往往只剩一两个片段，有时整个查询在一个场地就能完成。}
\Q{怎样判断片段和查询条件矛盾？}
\A{把片段条件和查询条件做合取，看是否恒假。例如片段 $E_{21}$ 的条件是 Bid=\q{A2}，查询要求 Bid=\q{A1}，合取为假，所以 $\sel{Bid=\q{A1}}(E_{21})=\varnothing$。答题时可以列一张表，写出片段、片段条件、查询条件和是否矛盾，作为中间过程。}
\Q{导出片段之间做连接时，怎样删分支？}
\A{用准则 2。$B_i$ 是在 $A_i$ 有贷款的企业，$L_{1j}$ 是 A1 的贷款。如果每个企业只在一个分行贷款，$L_{1j}\jn B_2$ 和 $L_{1j}\jn B_3$ 都为空，只保留 $L_{1j}\jn B_1$。}
\Q{什么情况下可以删垂直片段？}
\A{一个垂直片段除主键外的属性既不输出、也不出现在条件里，就可以删（准则 4）。例如只查 ENO、ENAME、BIRTH，这三个属性都在 $E_1$ 里，$E_2$ 的子树可以整个删去 \kj{5-68--69}。如果输出要 Sname、条件要 Dno，而这两个属性分在两个垂直片段里，两个片段都要保留，用 Sno 连接。}
\end{dlg}
\begin{pit}
\begin{itemize}
\item 题目要求写中间过程时，至少画三棵树：初始树、全局优化后的树、分片优化后的树。分片优化可以再加一棵本地化后的树或一张剪枝表。
\item 片段条件已经蕴含的选择可以去掉，例如 $\sel{Dno=2}(\R{Student_{12}})=\R{Student_{12}}$，保留也不算错。
\item 片段有副本时（如 $E_{11}$ 同时在 A1 和 A0），分片优化不需要选副本。选副本属于存取优化，可以补一句“选 A1 上的副本，本地执行”。
\end{itemize}
\end{pit}

% ---------------------------------------------------------------
\Needspace{14\baselineskip}
\subsection{2019 年第 1 题（20 分）}
\begin{prob}[原题]
某银行由总行和三个分行组成，三个分行编号分别为 $A_1$、$A_2$ 和 $A_3$，分别执行具体的放贷业务，总行没有放贷业务（编号 $A_0$）。现对贷款企业进行资产管理，涉及到的全局关系模式如下：\\
\textbf{Branchers}(Bid, Bname, Addrs, Mgrno)：分行基本信息表。Bid 为分行编号，Bname 为分行名称，Adrs 是分行地址，Mgrno 是分行负责人员工编号。\\
\textbf{Employee}(Eid, Ename, Etitle, Bid)：银行员工基本信息表。Eid 为员工编号，Ename 为员工姓名，Etitle 为员工职务，Bid 是员工所属的分行编号。\\
\textbf{Loanbusiness}(Lid, Lname, Lcapital, Ladr, Lpid)：贷款企业信息表。Lid 为企业编号，Lname 为企业名称，Lcapital 为企业注册资本，Ladr 为企业地址，Lpid 为企业负责人身份证号码。\\
\textbf{Loaninfo}(Bid, Lid, Eid, Lamount, Ldate, Lquota)：贷款信息表。Bid 为分行编号，Lid 为借贷企业编号，Eid 为银行贷款员编号，Lamount 为贷款金额，Ldate 为放贷日期，Lquota 为贷款期限。\\
数据管理的基本需求是：(1) 所有分行的基本信息都保存在总行场地；(2) 总行和各分行保存各自的员工信息，职务为“贷款员”(Etitle=\q{Loan operator}) 的员工信息在总行保存副本；(3) 各分行保存本分行员工发放的贷款信息，但单笔贷款金额 $\ge$1000 万元的贷款信息在总行保存副本；(4) 企业信息保存在发生借贷关系的分行。\\
1) 按要求进行分片设计，要求写出分片定义，画出分片树，并指出分片类型；\\
2) 按要求完成分配设计，指出分配类型；\\
3) 查询分行 $A_1$ 贷款员张龙向沈阳地区的企业发放的单笔贷款金额在 500 万元及以下的贷款信息。要求查询结果包括贷款企业名称、企业地址、贷款金额、放贷日、贷款期限。请写出查询的关系代数形式和 SQL 语句形式，并画出全局查询树；\\
4) 对 3) 题的查询进行全局优化，要求写出中间过程；\\
5) 在全局优化的基础上进行片段优化，要求写出中间过程。
\end{prob}

\begin{dlg}
\Q{四条需求分别涉及哪个关系？}
\A{(1) 只涉及 Branchers，整体放在 $A_0$，不分片。(2) 涉及 Employee：按 Bid 水平分片，为了副本再按是否为贷款员切一次。(3) 涉及 Loaninfo：按 Bid 水平分片，为了副本再按 Lamount$\ge$1000 切一次。(4) 涉及 Loanbusiness：划分依据“与哪个分行借贷”在 Loaninfo 里，所以是导出水平分片。}
\Q{总行员工 $E_0$ 要不要再按贷款员切？}
\A{不用。$E_0$ 本来就存在 $A_0$，“在总行存副本”对它没有意义；总行也没有放贷业务，没有贷款员。}
\Q{一个企业同时在 A1 和 A2 贷款，导出片段 $B_1$、$B_2$ 不就相交了吗？}
\A{会相交，这时按需求它确实要存在两个分行，相当于有副本。答卷里写一个假设：“设每个企业只与一个分行发生借贷关系，则 $B_1,B_2,B_3$ 两两不相交”。另外，没贷过款的企业不在任何 $B_i$ 里，所以还要假设每个企业至少有一笔贷款，才满足完备性。}
\end{dlg}

\begin{ans}[1) 分片定义、分片树、分片类型]
\textbf{Branchers}：不分片，整体存放在 $A_0$。\\
\textbf{Employee}：水平分片，分片属性 Bid、Etitle，记 LO = \q{Loan operator}。\quad
$E_0=\sel{Bid=\q{A0}}(\R{Employee})$；对 $i=1,2,3$：\\
\hspace*{2em}$E_{i1}=\sel{Bid=\q{Ai}\mand Etitle=LO}(\R{Employee})$，\quad
$E_{i2}=\sel{Bid=\q{Ai}\mand Etitle$\neq$LO}(\R{Employee})$。\\
\textbf{Loaninfo}：水平分片，分片属性 Bid、Lamount（单位万元）。对 $i=1,2,3$：\\
\hspace*{2em}$L_{i1}=\sel{Bid=\q{Ai}\mand Lamount$\ge$1000}(\R{Loaninfo})$，\quad
$L_{i2}=\sel{Bid=\q{Ai}\mand Lamount$<$1000}(\R{Loaninfo})$。\\
\textbf{Loanbusiness}：由 Loaninfo 的 Bid 导出的水平分片。对 $i=1,2,3$：\\
\hspace*{2em}$B_i=\R{Loanbusiness}\sjon{Lid}\sel{Bid=\q{Ai}}(\R{Loaninfo})=\R{Loanbusiness}\sjon{Lid}(L_{i1}\cup L_{i2})$。\\
假设每个企业至少有一笔贷款，且只与一个分行有借贷；Bid 只取 A0–A3。各层条件互斥且覆盖全部取值，满足完备性、可重构性（用 $\cup$ 重构）和不相交性。\\[2pt]
\centerline{%
\begin{forest} ft
[Employee [$E_0$,h] [$E_1$,h [$E_{11}$,h][$E_{12}$,h]] [$E_2$,h [$E_{21}$,h][$E_{22}$,h]] [$E_3$,h [$E_{31}$,h][$E_{32}$,h]]]
\end{forest}\quad
\begin{forest} ft
[Loaninfo [$L_1$,h [$L_{11}$,h][$L_{12}$,h]] [$L_2$,h [$L_{21}$,h][$L_{22}$,h]] [$L_3$,h [$L_{31}$,h][$L_{32}$,h]]]
\end{forest}\quad
\begin{forest} ft
[Loanbusiness [$B_1$,h][$B_2$,h][$B_3$,h]]
\end{forest}\quad
\begin{forest} ft
[Branchers]
\end{forest}}
\end{ans}

\begin{ans}[2) 分配设计与分配类型]
{\small\begin{tabular}{@{}ll@{}}
$A_0$（总行）：& Branchers，$E_0$，$E_{11},E_{21},E_{31}$（副本），$L_{11},L_{21},L_{31}$（副本）\\
$A_i$（$i=1,2,3$）：& $E_{i1},E_{i2}$，$L_{i1},L_{i2}$，$B_i$
\end{tabular}}\\
$E_{i1}$、$L_{i1}$ 存放在两个场地，其余片段只存放在一个场地，是部分复制分配。
\end{ans}

\begin{ans}[3) SQL、关系代数、全局查询树]
\begin{minipage}[t]{0.5\linewidth}\small\vspace{0pt}
\texttt{SELECT Lname, Ladr, Lamount, Ldate, Lquota}\\
\texttt{FROM\ \ \,Employee E, Loaninfo L, Loanbusiness B}\\
\texttt{WHERE E.Eid=L.Eid AND L.Lid=B.Lid}\\
\texttt{\ \ AND E.Ename='}张龙\texttt{' AND E.Bid='A1'}\\
\texttt{\ \ AND E.Etitle='Loan operator'}\\
\texttt{\ \ AND L.Bid='A1' AND L.Lamount<=500}\\
\texttt{\ \ AND B.Ladr='}沈阳\texttt{';}
\end{minipage}\hfill
\begin{minipage}[t]{0.46\linewidth}\centering\vspace{0pt}
\begin{forest} qt
[{$\proj{Lname,Ladr,Lamount,Ldate,Lquota}$}
 [{$\sel{$F_E\wedge F_L\wedge F_B$}$}
  [{$\jnon{L.Lid=B.Lid}$}
   [{$\jnon{E.Eid=L.Eid}$}[Employee][Loaninfo]]
   [Loanbusiness]]]]
\end{forest}
\end{minipage}\par\vspace{3pt}
关系代数：$\proj{Lname,Ladr,Lamount,Ldate,Lquota}\big(\sel{$F_E\wedge F_L\wedge F_B$}((\R{Employee}\jnon{Eid}\R{Loaninfo})\jnon{Lid}\R{Loanbusiness})\big)$，其中\\
\hspace*{1em}$F_E$：Ename=\q{张龙} $\wedge$ Etitle=\q{Loan operator} $\wedge$ Bid=\q{A1}\\
\hspace*{1em}$F_L$：Bid=\q{A1} $\wedge$ Lamount $\le$ 500\\
\hspace*{1em}$F_B$：Ladr=\q{沈阳}
\end{ans}
\begin{pit}
Employee 和 Loaninfo 都有 Bid，写成自然连接会把 Bid 也当作连接属性，所以要写带条件的连接 $\jnon{Eid}$。条件 L.Bid=\q{A1} 来自题意（A1 发放的贷款），写出来以后，分片优化能删掉 A2、A3 的贷款片段。
\end{pit}

\begin{ans}[4) 全局优化（中间过程）]
步骤 1（选择串接律）：$\sel{$F_E\wedge F_L\wedge F_B$}=\sel{$F_E$}\,\sel{$F_L$}\,\sel{$F_B$}$。\\
步骤 2（选择与连接的分配律）：$F_E$ 只含 Employee 的属性，移到 Employee 上方；$F_L$ 移到 Loaninfo 上方；$F_B$ 移到 Loanbusiness 上方。\\
步骤 3（投影串接律、投影与连接的分配律）：每个叶子上方只保留输出属性和连接属性。Employee 留 \{Eid\}；Loaninfo 留 \{Eid, Lid, Lamount, Ldate, Lquota\}；Loanbusiness 留 \{Lid, Lname, Ladr\}；第一次连接后留 \{Lid, Lamount, Ldate, Lquota\}。\\[2pt]
\centerline{%
\begin{forest} qt
[{$\proj{Lname,Ladr,Lamount,Ldate,Lquota}$}
 [{$\jnon{Lid}$}
  [{$\proj{Lid,Lamount,Ldate,Lquota}$}
   [{$\jnon{Eid}$}
    [{$\proj{Eid}$}[{$\sel{$F_E$}$}[Employee]]]
    [{$\proj{Eid,Lid,Lamount,Ldate,Lquota}$}[{$\sel{$F_L$}$}[Loaninfo]]]]]
  [{$\proj{Lid,Lname,Ladr}$}[{$\sel{$F_B$}$}[Loanbusiness]]]]]
\end{forest}}
\end{ans}

\begin{ans}[5) 分片优化（中间过程）]
步骤 1（数据本地化）：Employee $\to\cup(E_0,E_{11},E_{12},E_{21},E_{22},E_{31},E_{32})$；Loaninfo $\to\cup(L_{11},\dots,L_{32})$；Loanbusiness $\to\cup(B_1,B_2,B_3)$。\\
步骤 2（准则 1：选择作用于片段，矛盾的置空）：\\[1pt]
{\small\begin{tabular}{@{}llll@{}}
\toprule
片段 & 片段条件 & 与查询条件合取 & 结果\\ \midrule
$E_0$，$E_{2*}$，$E_{3*}$ & Bid=A0 / A2 / A3 & 与 Bid=A1 矛盾 & 空，删去 \\
$E_{12}$ & Bid=A1 $\wedge$ Etitle$\neq$LO & 与 Etitle=LO 矛盾 & 空，删去 \\
$E_{11}$ & Bid=A1 $\wedge$ Etitle=LO & 相容 & 保留 \\
$L_{11}$ & Bid=A1 $\wedge$ Lamount$\ge$1000 & 与 Lamount$\le$500 矛盾 & 空，删去\\
$L_{2*}$，$L_{3*}$ & Bid=A2 / A3 & 与 Bid=A1 矛盾 & 空，删去\\
$L_{12}$ & Bid=A1 $\wedge$ Lamount$<$1000 & 相容 & 保留\\
\bottomrule
\end{tabular}}\\[2pt]
步骤 3（准则 3、准则 2）：把 $\jnon{Lid}$ 移到 $\cup(B_1,B_2,B_3)$ 下面，得到 $L_{12}\jn B_1$、$L_{12}\jn B_2$、$L_{12}\jn B_3$。$B_2,B_3$ 中的企业不与 A1 借贷，连接条件不可能满足，所以 $L_{12}\jn B_2=L_{12}\jn B_3=\varnothing$，只保留 $B_1$。\\
步骤 4：去掉片段条件已经蕴含的谓词，得到最终的片段查询树。$E_{11}$（选 A1 上的副本）、$L_{12}$、$B_1$ 都在 $A_1$，查询在 $A_1$ 本地完成。\\[2pt]
\centerline{%
\begin{forest} qt
[{$\proj{Lname,Ladr,Lamount,Ldate,Lquota}$}
 [{$\jnon{Lid}$}
  [{$\proj{Lid,Lamount,Ldate,Lquota}$}
   [{$\jnon{Eid}$}
    [{$\proj{Eid}$}[{$\sel{Ename=\q{张龙}}$}[$E_{11}$]]]
    [{$\proj{Eid,Lid,Lamount,Ldate,Lquota}$}[{$\sel{Lamount$\le$500}$}[$L_{12}$]]]]]
  [{$\proj{Lid,Lname,Ladr}$}[{$\sel{Ladr=\q{沈阳}}$}[$B_1$]]]]]
\end{forest}}
\end{ans}
\begin{dlg}
\Q{不做“每个企业只在一个分行贷款”的假设，只留 $B_1$ 还对吗？}
\A{也对。$L_{12}$ 里的贷款都是 A1 发放的，借款企业都与 A1 有借贷关系，一定在 $B_1$ 里。$L_{12}\jn B_2$ 即使不为空，查出的也只是 $B_1$ 里已有企业的重复元组，所以结果不变。}
\end{dlg}

% ---------------------------------------------------------------
\Needspace{14\baselineskip}
\subsection{2020 年第 1 题（20 分）}
\begin{prob}[原题]
某机构由总部 ($A_0$) 和两个下属分支机构 ($A_1$ 和 $A_2$) 组成。总部和分支机构都从事科研业务。现对科研项目信息进行数据管理，涉及到的全局关系模式如下：\\
\textbf{Branchers}(Bid, Bname, Addrs, Mgrno)：机构基本信息表。Bid 为机构编号，Bname 为机构名称，Adrs 是机构地址，Mgrno 是机构负责人员工编号。\\
\textbf{Employee}(Eid, Ename, Etitle, Bid)：员工基本信息表。Eid 为员工编号，Ename 为员工姓名，Etitle 为员工职务，Bid 是员工所属的机构编号。\\
\textbf{Projects}(Pid, Pname, Budget, DESC, Security, Pleader)：项目基本信息表。Pid 为项目编号，Budget 为项目金额，DESC 为项目描述信息，Security 为项目是否涉密，Pleader 为项目负责人员工编号。\\
数据管理的基本需求是：(1) 非涉密项目保存在项目负责人所在的分支机构场地，其中金额超过 100 万元（含）的项目信息在总行保存副本；(2) 涉密项目信息只在总部场地保存；(3) 总部和各分支机构保存各自的员工信息；(4) 涉密项目的负责人信息在总行保存副本；(5) 机构信息仅在总部场地保存。\\
请回答：1) 按要求对表进行分片。写出分片定义，画出分片树，并指出分片类型；2) 按需求完成分配设计，指出分配类型；3) 查询机构 $A_2$ 员工张龙负责的项目金额不超过 80 万的非涉密项目信息。要求查询结果包括张龙的姓名、项目名称、项目描述信息和项目金额；4) 对 3) 题的查询进行全局优化，要求写出中间过程；5) 在全局优化的基础上进行片段优化，要求写出中间过程。
\end{prob}
\begin{dlg}
\Q{这题和 2019 年那题有什么不同？}
\A{这题有两处导出。一是项目存在“负责人所在场地”，Projects 里只有 Pleader，场地要到 Employee.Bid 里查，所以 Projects 由 Employee 导出。二是“涉密项目的负责人”要在总部存副本，一个员工是不是涉密项目负责人要到 Projects 里查，所以 Employee 也要由 Projects 导出，再切一次。}
\Q{两个关系互相导出，会不会循环定义？}
\A{不会，只要都用全局关系来定义。例如 $E_{i1}=\sel{Bid=\q{Ai}}(\R{Employee})\sjon{Eid=Pleader}\sel{Security=\q{涉密}}(\R{Projects})$，右边都是全局关系，不依赖其他片段。}
\Q{总部员工负责的非涉密项目放在哪里？}
\A{题目说的是“分支机构场地”，所以默认非涉密项目的负责人都在 A1、A2，答卷里写一句假设。如果总部员工也可能负责非涉密项目，就再加 $P_{01},P_{02}$ 放在 A0。}
\end{dlg}

\begin{ans}[1) 分片定义、分片树、分片类型]
\textbf{Branchers}：不分片，存放在 $A_0$。\\
\textbf{Projects}：水平分片（Security、Budget 上的基本水平分片，加上按负责人 Bid 的导出水平分片）。\\
\hspace*{1em}$P_S=\sel{Security=\q{涉密}}(\R{Projects})$；\quad 对 $i=1,2$：\\
\hspace*{1em}$P_{i1}=\sel{Security=\q{非涉密}\mand Budget$\ge$100}\big(\R{Projects}\sjon{Pleader=Eid}\sel{Bid=\q{Ai}}(\R{Employee})\big)$，\\
\hspace*{1em}$P_{i2}=\sel{Security=\q{非涉密}\mand Budget$<$100}\big(\R{Projects}\sjon{Pleader=Eid}\sel{Bid=\q{Ai}}(\R{Employee})\big)$。\\
\textbf{Employee}：水平分片（Bid 上的基本水平分片，加上由 Projects 导出的分片）。$E_0=\sel{Bid=\q{A0}}(\R{Employee})$；对 $i=1,2$：\\
\hspace*{1em}$E_{i1}=\sel{Bid=\q{Ai}}(\R{Employee})\sjon{Eid=Pleader}\sel{Security=\q{涉密}}(\R{Projects})$（A$i$ 中的涉密项目负责人），\\
\hspace*{1em}$E_{i2}=\sel{Bid=\q{Ai}}(\R{Employee})-E_{i1}$（A$i$ 的其他员工）。\\
假设非涉密项目的负责人都属于 A1 或 A2；Bid 只取 A0–A2。\\[2pt]
\centerline{%
\begin{forest} ft
[Projects [$P_S$,h] [$P_N$,h [$P_{N1}$,h [$P_{11}$,h][$P_{12}$,h]] [$P_{N2}$,h [$P_{21}$,h][$P_{22}$,h]]]]
\end{forest}\qquad
\begin{forest} ft
[Employee [$E_0$,h] [$E_1$,h [$E_{11}$,h][$E_{12}$,h]] [$E_2$,h [$E_{21}$,h][$E_{22}$,h]]]
\end{forest}\qquad
\begin{forest} ft
[Branchers]
\end{forest}}
{\small $P_N$ 是非涉密项目，$P_{Ni}$ 是负责人在 A$i$ 的非涉密项目。}
\end{ans}

\begin{ans}[2) 分配设计]
{\small\begin{tabular}{@{}ll@{}}
$A_0$（总部）：& Branchers，$E_0$，$P_S$，$E_{11},E_{21}$（副本），$P_{11},P_{21}$（副本）\\
$A_1$：& $E_{11},E_{12},P_{11},P_{12}$\\
$A_2$：& $E_{21},E_{22},P_{21},P_{22}$
\end{tabular}}\\
部分片段有两个副本，其余只有一个，是部分复制分配。
\end{ans}

\begin{ans}[3) SQL、关系代数、全局查询树]
\begin{minipage}[t]{0.52\linewidth}\small\vspace{0pt}
\texttt{SELECT E.Ename, P.Pname, P.DESC, P.Budget}\\
\texttt{FROM\ \ \,Employee E, Projects P}\\
\texttt{WHERE E.Eid=P.Pleader AND E.Ename='}张龙\texttt{'}\\
\texttt{\ \ AND E.Bid='A2' AND P.Budget<=80}\\
\texttt{\ \ AND P.Security='}非涉密\texttt{';}
\end{minipage}\hfill
\begin{minipage}[t]{0.44\linewidth}\centering\vspace{0pt}
\begin{forest} qt
[{$\proj{Ename,Pname,DESC,Budget}$}
 [{$\sel{$F_E\wedge F_P$}$}
  [{$\jnon{Eid=Pleader}$}[Employee][Projects]]]]
\end{forest}
\end{minipage}\par\vspace{3pt}
关系代数：$\proj{Ename,Pname,DESC,Budget}\big(\sel{$F_E\wedge F_P$}(\R{Employee}\jnon{Eid=Pleader}\R{Projects})\big)$，\\
\hspace*{1em}$F_E$：Ename=\q{张龙}$\wedge$Bid=\q{A2}；\quad $F_P$：Budget$\le$80$\wedge$Security=\q{非涉密}。
\end{ans}

\begin{ans}[4) 全局优化 \quad 5) 分片优化]
4) 选择串接律拆成 $\sel{$F_E$}\sel{$F_P$}$；选择与连接的分配律把 $F_E$ 移到 Employee 上方，$F_P$ 移到 Projects 上方；用投影的分配律后，Employee 留 \{Eid, Ename\}，Projects 留 \{Pleader, Pname, DESC, Budget\}（左树）。\\
5) 本地化：Employee $\to\cup(E_0,E_{11},E_{12},E_{21},E_{22})$，Projects $\to\cup(P_S,P_{11},P_{12},P_{21},P_{22})$。\\
准则 1：Bid=\q{A2} 与 $E_0,E_{11},E_{12}$ 矛盾，删去，留 $E_{21},E_{22}$（不知道张龙是否为涉密项目负责人，两个片段都留）；
Security=\q{非涉密} 与 $P_S$ 矛盾，Budget$\le$80 与 $P_{11},P_{21}$（Budget$\ge$100）矛盾，删去，留 $P_{12},P_{22}$。\\
准则 3、准则 2：连接移到并下面，得到 4 个连接。$P_{12}$ 的负责人属于 A1，与 A2 的员工按 Eid=Pleader 连接必为空，所以 $E_{2j}\jn P_{12}=\varnothing$，只剩 $E_{21}\jn P_{22}$、$E_{22}\jn P_{22}$。这些片段都在 $A_2$，本地执行（右树）。\\[2pt]
\begin{forest} qt
[{$\proj{Ename,Pname,DESC,Budget}$}
 [{$\jnon{Eid=Pleader}$}
  [{$\proj{Eid,Ename}$}[{$\sel{$F_E$}$}[Employee]]]
  [{$\proj{Pleader,Pname,DESC,Budget}$}[{$\sel{$F_P$}$}[Projects]]]]]
\end{forest}\hfill
\begin{forest} qt
[{$\proj{Ename,Pname,DESC,Budget}$}
 [{$\cup$}
  [{$\jnon{Eid=Pleader}$}
   [{$\proj{Eid,Ename}$}[{$\sel{Ename=\q{张龙}}$}[$E_{21}$]]]
   [{$\proj{Pleader,\dots}$}[{$\sel{Budget$\le$80}$}[$P_{22}$]]]]
  [{$\jnon{Eid=Pleader}$}
   [{$\proj{Eid,Ename}$}[{$\sel{Ename=\q{张龙}}$}[$E_{22}$]]]
   [{$\proj{Pleader,\dots}$}[{$\sel{Budget$\le$80}$}[$P_{22}$]]]]]]
\end{forest}
\end{ans}

% ---------------------------------------------------------------
\Needspace{14\baselineskip}
\subsection{2023 年第一题}
\begin{prob}[原题]
下面是某高校数据库的全局模式：学生 Student = \{Sno, Sname, HomeAddr, Phone, Class, Birth, Dno\}；学院 DEPT = \{Dno, Dname\}。该学校共有 3 个学院，Dno 分别为 0，1，2。要求 DEPT 关系和 Student 关系的部分属性 \{Sno, Sname, HomeAddr, Phone\} 保存在学校场地 S0 上，Student 关系的那部分属性 \{Sno, Class, Birth, Dno\} 保存在学生所在学院场地 S1、S2 和 S3 上。要求：\\
1. 将全局模式进行分片，写出分片定义和分片条件；\\
2. 指出各分片的类型（垂直、水平分片、混合），并画出分片树；\\
3. 对查询 \texttt{SELECT Sname, Class, Dname FROM Student, DEPT WHERE Dno = 2 AND Student.Dno = DEPT.Dno;} 进行全局优化，画出优化后的全局查询树，要求写出中间过程；\\
4. 进行分片优化，画出优化后的分片查询树，要求写出中间过程。
\end{prob}
\begin{ans}[1. 分片定义 \quad 2. 类型与分片树]
\begin{minipage}[t]{0.6\linewidth}
DEPT：不分片，存放在 S0。\\
Student：先垂直、后水平，是混合分片。\\
\hspace*{1em}$\R{Student_0}=\proj{Sno,Sname,HomeAddr,Phone}(\R{Student})$（存 S0）\\
\hspace*{1em}$\R{Student_1}=\proj{Sno,Class,Birth,Dno}(\R{Student})$\\
\hspace*{1em}$\R{Student_{1j}}=\sel{Dno=$j$}(\R{Student_1})$，$j=0,1,2$（依次存 S1、S2、S3）\\
验证：$\text{Attr}(S_0)\cup\text{Attr}(S_1)=\text{Attr}(\R{Student})$，$\text{Attr}(S_0)\cap\text{Attr}(S_1)=\{\text{Sno}\}$，即主键；三个水平条件互斥且覆盖。\\
分配：S0 存 DEPT、Student$_0$；S1 存 Student$_{10}$；S2 存 Student$_{11}$；S3 存 Student$_{12}$。没有片段被复制，是全分割。
\end{minipage}\hfill
\begin{minipage}[t]{0.36\linewidth}\centering\vspace{2pt}
\begin{forest} ft
[Student [Student$_0$,v][Student$_1$,v [S$_{10}$,h][S$_{11}$,h][S$_{12}$,h]]]
\end{forest}\quad
\begin{forest} ft
[DEPT]
\end{forest}
\end{minipage}
\end{ans}

\begin{ans}[3. 全局优化 \quad 4. 分片优化]
初始查询：$\proj{Sname,Class,Dname}(\sel{Dno=2}(\R{Student}\jnon{Dno}\R{DEPT}))$（图 a）。\\
全局优化：Dno=2 通过连接条件传递到 DEPT，两边都下移；Student 上方投影 \{Sname, Class, Dno\}，DEPT 上方投影 \{Dno, Dname\}（图 b）。\\
分片优化：\cir{1} 本地化，Student $\to$ Student$_0\jnon{Sno}\cup$(Student$_{10}$, Student$_{11}$, Student$_{12}$)（图 c）；
\cir{2} Dno 只在垂直片段 Student$_1$ 上，$\sel{Dno=2}$ 移进 $\cup$ 后与 Student$_{10}$、Student$_{11}$ 矛盾，删去，只留 Student$_{12}$，而且 $\sel{Dno=2}(\R{Student_{12}})=\R{Student_{12}}$；
\cir{3} 投影分到两个垂直片段：Student$_0$ 留 \{Sno, Sname\}，Student$_{12}$ 留 \{Sno, Class, Dno\}，保留 Sno 用于垂直重构的连接；HomeAddr、Phone、Birth 用不到。两个垂直片段都提供输出属性，都不能删（图 d）。\\[3pt]
\begin{tabular}{@{}c@{\hspace{6pt}}c@{\hspace{6pt}}c@{\hspace{6pt}}c@{}}
\begin{forest} qt
[{$\proj{Sname,Class,Dname}$}[{$\sel{Dno=2}$}[{$\jnon{Dno}$}[Student][DEPT]]]]
\end{forest} &
\begin{forest} qt
[{$\proj{Sname,Class,Dname}$}
 [{$\jnon{Dno}$}
  [{$\proj{Sname,Class,Dno}$}[{$\sel{Dno=2}$}[Student]]]
  [{$\proj{Dno,Dname}$}[{$\sel{Dno=2}$}[DEPT]]]]]
\end{forest} &
\begin{forest} qt
[{$\proj{Sname,Class,Dname}$}
 [{$\jnon{Dno}$}
  [{$\proj{Sname,Class,Dno}$}[{$\sel{Dno=2}$}
    [{$\jnon{Sno}$}[S$_0$][{$\cup$}[S$_{10}$,cut][S$_{11}$,cut][S$_{12}$]]]]]
  [{$\proj{Dno,Dname}$}[{$\sel{Dno=2}$}[DEPT]]]]]
\end{forest} &
\begin{forest} qt
[{$\proj{Sname,Class,Dname}$}
 [{$\jnon{Dno}$}
  [{$\proj{Sname,Class,Dno}$}
   [{$\jnon{Sno}$}[{$\proj{Sno,Sname}$}[S$_0$]][{$\proj{Sno,Class,Dno}$}[S$_{12}$]]]]
  [{$\proj{Dno,Dname}$}[{$\sel{Dno=2}$}[DEPT]]]]]
\end{forest}\\
\small (a) 初始 & \small (b) 全局优化后 & \small (c) 本地化，灰色已删去 & \small (d) 分片优化后
\end{tabular}
\end{ans}

\clearpage
\section{代价计算：I/O 代价、执行场地、半连接}

\subsection{I/O 代价（2023 年第三题，课件 5-11--18 的原例）}\label{s:io}
\begin{prob}[原题]
假设有学生关系 Students\{Sno, Sname, Class\}，选课关系 SC\{Sno, Cno\}，选修 5 号课程的学生姓名 SQL 语句如下：\\
\texttt{SELECT Sname FROM Student, SC WHERE Student.Sno = SC.Sno AND Cno = '5';}\\
可以转换成三种关系代数表达式：\\
\hspace*{1em}$Q_1=\proj{Sname}(\sel{Student.Sno=SC.Sno$\wedge$Cno=5}(\R{Student}\times\R{SC}))$\\
\hspace*{1em}$Q_2=\proj{Sname}(\sel{Cno=5}(\R{Student}\jnon{Sno}\R{SC}))$，假设自然连接结果大小为 $10^4$ 个元组\\
\hspace*{1em}$Q_3=\proj{Sname}(\R{Student}\jnon{Sno}\sel{Cno=5}(\R{SC}))$\\
假设内存中存放 5 块 Students 元组和一块 SC 元组，一块可以装 10 个 Students 元组或 100 个 SC 元组。Students 有 1000 个元组，SC 有 10000 个元组，其中选 5 号课程的有 50 个元组，采用循环全连接算法实现，假设 I/O 读写速度是每秒 25 块。计算 $Q_1$、$Q_2$ 和 $Q_3$ 的 I/O 读写时间代价。
\end{prob}
\begin{flow}[计算步骤]
\begin{enumerate}
\item 块数 = 元组数 ÷ 每块元组数：Students $1000/10=100$ 块，SC $10000/100=100$ 块。
\item 嵌套循环连接读的块数 = 外层块数 + 外层批数 × 内层块数，外层批数 = 外层块数 ÷ 外层可用的内存块数。
\item 中间结果放不下内存，就要先写到外存，下一步再读回来，代价是 2 × 中间结果块数。结果能放进内存（如 50 个元组）就不计 I/O。
\item 时间 = 块数 ÷ 每秒块数，各步相加。
\end{enumerate}
\end{flow}
\begin{dlg}
\Q{“SC 被遍历的次数 $=1000/(10\times5)=20$”是怎么来的？}
\A{内存一次只能放 5 块 Students，也就是 50 个学生。Students 共 100 块，要分 $100/5=20$ 批读入。每读入一批，都要把 SC 从头到尾逐块读一遍（每次读进那 1 块 SC 缓冲区），和这 50 个学生配对。所以 SC 被读 20 遍，每遍 100 块。}
\Q{$Q_1$ 的笛卡尔积为什么要写出去再读回来？}
\A{$1000\times10^4=10^7$ 个元组放不进内存，只能先写到磁盘（$10^7/10=10^6$ 块），做选择时再读回来，又是 $10^6$ 块，所以算两次。每块装 10 个笛卡尔积元组是课件和参考答案用的假设，答题时要写出来。}
\Q{$Q_3$ 为什么只要 8 秒？}
\A{它先做选择：读一遍 SC（100 块），筛完只剩 50 个元组，全部放进内存。连接时把 Students 读一遍（100 块），逐个和内存里的 50 个元组配对，不需要反复扫描，也不写中间结果。}
\end{dlg}
\begin{ans}[2023 年第三题]
假设：笛卡尔积和连接的中间结果每块可装 10 个元组。\\
$Q_1$：\cir{1} 读：$\underbrace{1000/10}_{\text{Students 块数}}+\underbrace{1000/(10\times5)}_{\text{SC 遍历次数}}\times\underbrace{10000/100}_{\text{SC 块数}}=100+20\times100=2100$ 块，$2100/25=84$\,s；\\
\cir{2} 写笛卡尔积：$10^3\times10^4=10^7$ 个元组，$10^7/10=10^6$ 块，$10^6/25=4\times10^4$\,s；\cir{3} 选择时读回：$4\times10^4$\,s，选择后剩 50 个元组，放在内存；\cir{4} 投影不需要 I/O。\\
\hspace*{1em}共 $84+2\times4\times10^4=80084\,\text{s}$，约 22.2 小时。\\
$Q_2$：\cir{1} 读：同上，2100 块，84\,s；\cir{2} 写连接结果：$10^4/10=1000$ 块，40\,s；\cir{3} 选择时读回：40\,s；\cir{4} 投影不需要 I/O。共 $84+40+40=164\,\text{s}$，约 2.7 分钟。\\
$Q_3$：\cir{1} 选择：读 SC $10000/100=100$ 块，4\,s，结果 50 个元组在内存；\cir{2} 连接：读 Students $1000/10=100$ 块，4\,s，结果 50 个元组在内存；\cir{3} 投影不需要 I/O。共 8\,s。\\
结论：$Q_3$ 代价最小。它先选择、再连接，没有笛卡尔积，比 $Q_1$ 快了约一万倍。
\end{ans}
\begin{pit}
课件原例是“每秒读写 20 块、选 2 号课程”，结果分别是 105\,s、$10^5$\,s（约 28 小时）、205\,s、10\,s。考试换了数字，算法相同：先算块数，再除以速度。
\end{pit}

\subsection{选择执行场地（2023 年第四题，课件 6-5--7 的原例）}\label{s:site}
\begin{prob}[原题]
有关系 EMP 和 DEPT 如下：EMP\{ENO, ENAME, BIRTH, SALARY, DNO\}，DEPT\{DNO, DNAME\}。假设 EMP：元组数 10000 条，元组大小 100 Byte；DEPT：元组数 100 条，元组大小 35 Byte。存在三个场地 S1、S2 和 S3。EMP 存放在 S1 场地，DEPT 存放在 S2 场地，S3 为查询场地。查询每个雇员的姓名及所在单位名称 SQL 语句：\texttt{SELECT ENAME, DNAME FROM EMP, DEPT WHERE EMP.DNO = DEPT.DNO;} 关系表达式为 $\proj{ENAME,DNAME}(\R{EMP}\jnon{DNO}\R{DEPT})$。假设结果元组大小 40 字节，S3 为查询场地，结果关系大小：$40\times10000=400$\,KB。请选择传输代价最小的查询执行场地。
\end{prob}
\begin{flow}[计算步骤]
\begin{enumerate}
\item 算每个关系的大小：Size = Card × Length；结果大小 = 结果元组数 × 结果元组长度。
\item 对每个候选执行场地 $S_k$：传输量 = 不在 $S_k$ 的输入关系大小之和，若 $S_k$ 不是查询场地，再加上结果大小。
\item 取最小的。题目给了 $C_0,C_1$ 时，每次传输的代价是 $C_0+C_1\cdot X$。
\end{enumerate}
\end{flow}
\begin{ans}[2023 年第四题]
Size(EMP) $=10000\times100\,\text{B}=1000$\,KB；Size(DEPT) $=100\times35\,\text{B}=3.5$\,KB；Size(R) $=10000\times40\,\text{B}=400$\,KB。\\
\begin{tabular}{@{}lll@{}}
在 S1 执行：& 传 DEPT 到 S1，结果 R 传到 S3 & $3.5+400=403.5$\,KB\\
在 S2 执行：& 传 EMP 到 S2，结果 R 传到 S3 & $1000+400=1400$\,KB\\
在 S3 执行：& 传 EMP、DEPT 到 S3 & $1000+3.5=1003.5$\,KB
\end{tabular}\\
选 S1 为执行场地，传输代价最小，为 403.5\,KB。
\end{ans}
\begin{dlg}
\Q{在查询场地 S3 执行看起来最自然，为什么反而不好？}
\A{EMP 有 1000\,KB，整个搬走代价很高。一般把小关系搬到大关系所在的场地，但结果还要送回查询场地，所以结果大小也要算进去再比较。这题结果只有 400\,KB，比 EMP 小，所以在 S1 执行更省。}
\end{dlg}

\Needspace{14\baselineskip}
\subsection{直接连接与半连接的通信代价（2020、2019 年第 2 题，课件 6-15--39）}\label{s:semi}
\begin{flow}[估算与比较的步骤]
\begin{enumerate}
\item 写假设：数据均匀分布、属性相互独立；外键的取值来自被参照关系的主键；只计传输的数据量，或写成 $C_0+C_1X$。
\item 先做本地选择，不需要通信。等值条件 $\rho=1/\text{Val}(R,A)$；范围条件 $\rho=(X-\text{Min})/(\text{Max}-\text{Min})$，也可以按取值个数算；Card $=\rho\cdot$Card$(R)$。
\item 直接连接（全连接）：把一个关系整个传到另一个关系所在的场地做连接，结果再传回查询场地。两个方向都算，取小的。
\item 半连接：\cir{1} 把 $S'$ 的连接列 $\proj{A}(S')$ 传到 $R$ 的场地，大小为 $\text{Val}(S',A)\times\text{Length}(A)$；
\cir{2} 计算 $R'=R\sj S'$，$\rho=\text{Val}(S',A)/\text{Val}(S,A)$，Card$(R')=\rho\cdot$Card$(R)$；
\cir{3} 把 $R'$ 传回去做连接。
\item 比较：当 $C_0/C_1+\text{Size}(\proj{A}S')+\text{Size}(R')\le\text{Size}(R)$ 时，半连接比直接传 $R$ 省 \kj{6-39}。
\end{enumerate}
\end{flow}
\begin{dlg}
\Q{半连接省在哪里？}
\A{直接连接要把整个 $R$ 传过去，而 $R$ 里大部分元组连不上。半连接先把对方的连接列（很窄）传过来，筛掉连不上的元组，然后只传能连上的 $R'$。连不上的元组越多，也就是选择度 $\rho$ 越小，省得越多。}
\Q{什么时候半连接反而不划算？}
\A{$\rho$ 接近 1（几乎都能连上），或者连接列本身很宽时。多传的 $\proj{A}(S')$ 和多出的一次 $C_0$ 换不来多少缩减。半连接还会增加本地处理代价，在局域网里不一定合算 \kj{6-63}。}
\Q{半连接的选择度为什么是 $\text{Val}(S',A)/\text{Val}(S,A)$？}
\A{$R$ 的外键 $A$ 均匀地取自 $S$ 的 Val$(S,A)$ 个值。$S$ 筛成 $S'$ 后只剩 Val$(S',A)$ 个值，$R$ 中外键落在这些值里的比例就是两者之比。课件 SDD-1 例子里的 $\rho_1=200/2000=0.1$ 就是这样算的 \kj{6-51}。}
\end{dlg}

\begin{keep}
\begin{prob}[2020 年第 2 题（10 分）原题]
科研项目表存储在科技处场地，教师信息表存储在人事处场地，两个场地通过互联网连接。现从人事处场地查询计算机学院教师承担的项目信息。假设数据满足均匀分布，关系概要图如下：\\[2pt]
{\footnotesize\setlength{\tabcolsep}{3pt}
\begin{tabular}{@{}l|cccc@{}}
\multicolumn{5}{@{}l}{Projects（单位：字节）Card=30000 项，Length=200}\\ \hline
 & PID & FID & Pname & Sdate \\ \hline
length & 15 & 8 & 100 & 8 \\
val & 30000 & 2500 & 30000 & 25000 \\
说明 & 项目编号 & 负责人编号 & 项目名称 & 立项时间
\end{tabular}\hfill
\begin{tabular}{@{}l|ccc@{}}
\multicolumn{4}{@{}l}{Faculties（单位：字节）Card=3000 人，Length=300}\\ \hline
 & FID & Fname & Institute \\ \hline
length & 8 & 20 & 8 \\
val & 3000 & 2800 & 20 \\
说明 & 教师编号 & 教师姓名 & 学院
\end{tabular}}\\[3pt]
计算：(1) 如果采用全连接优化算法，请给出查询执行计划，并计算通信代价。(2) 请给出优化的查询执行策略，并计算通信代价。
\end{prob}
\begin{ans}[2020 年第 2 题]
假设：均匀分布；Projects.FID 的取值来自 Faculties.FID；结果是项目信息（Projects 元组，200\,B）；只计传输字节数，每次传输另加 $C_0$。\\
本地选择：$F'=\sel{Institute=\q{计算机}}(\R{Faculties})$，Card$(F')=3000/20=150$，Val$(F',\text{FID})=150$。\\
选择度 $\rho=150/3000=0.05$，计算机学院教师承担的项目数 $=30000\times0.05=1500$。\\[2pt]
(1) 全连接：\\
\hspace*{1em}方案 A，在人事处连接：Projects 整个传到人事处，$30000\times200=6\,000\,000$\,B；\\
\hspace*{1em}方案 B，在科技处连接：$F'$ 传到科技处 $150\times300=45\,000$\,B，连接后结果传回人事处 $1500\times200=300\,000$\,B，共 $345\,000$\,B。\\
\hspace*{1em}采用方案 B，通信代价 $345\,000$\,B，即 $2C_0+345\,000\,C_1$。\\[2pt]
(2) 半连接策略：\\
\hspace*{1em}\cir{1} 人事处把 $\proj{FID}(F')$ 传到科技处：$150\times8=1\,200$\,B；\\
\hspace*{1em}\cir{2} 科技处计算 $P'=\R{Projects}\sjon{FID}\proj{FID}(F')$，Card$(P')=1500$；\\
\hspace*{1em}\cir{3} $P'$ 传回人事处：$1500\times200=300\,000$\,B；\cir{4} 人事处计算 $F'\jnon{FID}P'$ 得到结果。\\
\hspace*{1em}通信代价 $1\,200+300\,000=301\,200$\,B，即 $2C_0+301\,200\,C_1$，小于全连接的 345\,000\,B，是最优策略。
\end{ans}
\end{keep}

\begin{keep}
\begin{prob}[2019 年第 2 题（10 分）原题]
校办拟查询 40 岁及以下教师的基本信息和科研信息，以及他们第 6 学期的教学信息。已知：教师的基本信息（Faculties 表）位于人事处、科研项目信息（Projects 表）位于科技处、教学信息（Schedules）位于教务处。各场地通过互联网连接，数据满足均匀分布。关系概要图如下：\\[2pt]
{\footnotesize\setlength{\tabcolsep}{3pt}
\begin{tabular}{@{}l|ccc@{}}
\multicolumn{4}{@{}l}{Faculties：Card=3000 人，Length=300}\\ \hline
 & FID & Fname & Fage \\ \hline
length & 8 & 20 & 1 \\
val & 3000 & 2800 & 40 \\
说明 & 教师编号 & 教师姓名 & 年龄
\end{tabular}\hfill
\begin{tabular}{@{}l|cccc@{}}
\multicolumn{5}{@{}l}{Projects：Card=30000 项，Length=200}\\ \hline
 & PID & FID & Pname & Sdate \\ \hline
length & 15 & 8 & 100 & 8 \\
val & 30000 & 2500 & 30000 & 25000 \\
说明 & 项目编号 & 负责人编号 & 项目名称 & 立项时间
\end{tabular}\\[4pt]
\begin{tabular}{@{}l|cccc@{}}
\multicolumn{5}{@{}l}{Schedules：Card=120000 门次，Length=150}\\ \hline
 & CID & Cname & Semester & FID \\ \hline
length & 15 & 60 & 1 & 8 \\
val & 4000 & 4000 & 8 & 2500 \\
说明 & 课程编号 & 课程名称 & 开课学期 & 主讲教师
\end{tabular}\hfill\parbox[b]{6.5cm}{学校教师的年龄均匀分布在 [21, 60] 区间。}}\\[2pt]
给出优化的查询策略，并计算查询代价。
\end{prob}
\begin{ans}[2019 年第 2 题]
假设：均匀分布、属性相互独立；Projects.FID 与 Schedules.FID 取自 Faculties.FID；年龄取整数 21–60；只计传输字节；校办本地没有数据。\\
\cir{1} 本地选择（不需要通信）：\\
\hspace*{1em}人事处 $F'=\sel{Fage$\le$40}(\R{Faculties})$：21–40 共 20 个值，$\rho=20/40=1/2$（按公式 $(40-21)/(60-21)\approx0.49$，近似 1/2），Card $=1500$，Size $=1500\times300=450\,000$\,B；\\
\hspace*{1em}教务处 $S_6=\sel{Semester=6}(\R{Schedules})$：$\rho=1/8$，Card $=15\,000$。\\
\cir{2} 半连接：$\proj{FID}(F')$（$1500\times8=12\,000$\,B）分别传到科技处和教务处，共 $24\,000$\,B。\\
\hspace*{1em}科技处 $P'=\R{Projects}\sj\proj{FID}(F')$：$\rho=1500/3000=1/2$，Card $=15\,000$，Size $=3\,000\,000$\,B；\\
\hspace*{1em}教务处 $S'=S_6\sj\proj{FID}(F')$：Card $=15\,000\times1/2=7\,500$，Size $=7500\times150=1\,125\,000$\,B。\\
\cir{3} 把 $F',P',S'$ 传到校办：$450\,000+3\,000\,000+1\,125\,000=4\,575\,000$\,B，在校办计算 $F'\jnon{FID}P'\jnon{FID}S'$。\\
总通信代价：$24\,000+4\,575\,000=4\,599\,000\,\text{B}\approx4.6$\,MB，共 5 次传输，即 $5C_0+4\,599\,000\,C_1$。\\
对比：只做本地选择、不用半连接，要传 $450\,000+6\,000\,000+15000\times150=8\,700\,000$\,B；什么都不做，要传 $900\,000+6\,000\,000+18\,000\,000=24\,900\,000$\,B。两个半连接的收益分别是 $0.5\times6\,000\,000$\,B 和 $0.5\times2\,250\,000$\,B，都远大于各自 12\,000\,B 的代价，所以采用。
\end{ans}
\end{keep}
\begin{dlg}
\Q{为什么不先在某个场地把三个关系连好，再把结果传给校办？}
\A{先估结果大小。每个老师平均有 $30000/2500=12$ 个项目、$15000/2500=6$ 门第 6 学期的课，三表按 FID 连接得到的是项目和课程的组合，约 8–9 万行，每行 600 多字节，约 50\,MB，比三个缩减后的关系加起来的 4.6\,MB 大得多。结果比输入大时，应该把缩减后的输入送到查询场地再连接。}
\Q{如果题目的意思是两张结果表（教师和科研、教师和教学），策略要变吗？}
\A{不用变。还是把 $F',P',S'$ 送到校办，分别做两次连接，通信代价一样。}
\end{dlg}

\clearpage
\section{分布式事务：子事务、进程模型、2PC 与故障处理}

\subsection{分解子事务与进程代理模型 \kj{7-19--26}}\label{s:proc}
\begin{flow}[分解子事务、写进程代理模型的步骤]
\begin{enumerate}
\item 给每一步操作标上它访问的数据所在的场地。一个场地对应一个子事务，写清每个子事务做什么。
\item 先像集中式那样写出整个事务（begin\_transaction … Commit … end），再拆分，阅卷时容易看清。
\item 拆成根代理和子代理。根代理在源场地，负责 input、条件判断、\texttt{Create AGENTi}、\texttt{Send(AGENTi,}\,参数\texttt{)}，最后执行 \texttt{Commit}，由它发起两阶段提交。子代理负责 \texttt{Receive(ROOT,}\,参数\texttt{)}，执行本场地的 SQL，检查不通过就 \texttt{Abort}。
\item 题目要求画图时，画进程模型：协调者进程（根代理，含全局事务）和各代理者进程（子代理，各含一个子事务）。
\end{enumerate}
\end{flow}
\begin{dlg}
\Q{根代理和子代理是什么关系？}
\A{每个应用只有一个根代理，它负责全局事务的开始、提交和废弃。只有根代理能创建子代理。所有子事务都能成功提交时，根代理才决定全局提交，否则全局废弃 \kj{7-22}。在 2PC 里，根代理是协调者，子代理是参与者。}
\Q{进程模型和服务器模型怎么区分？}
\A{进程模型为每个子事务在它的场地上建一个代理者进程，另有一个协调者进程。服务器模型在每个场地只建一个服务器进程，轮流为多个全局事务的子事务服务，不同全局事务在同一场地上共用它 \kj{7-26}。}
\Q{子代理里的 Abort 能决定全局废弃吗？}
\A{不能。它只让本场地的子事务失败。在 2PC 的决定阶段，这个子代理会回答“准备废弃”，全局废弃由协调者决定。}
\end{dlg}

\begin{keep}
\begin{prob}[2004 年第 3 题]
画出分布式事务进程代理模型，并举实例说明执行过程。
\end{prob}
\begin{ans}[2004 年第 3 题]
\begin{minipage}[t]{0.36\linewidth}\vspace{0pt}\centering
\begin{tikzpicture}[font=\footnotesize]
\node[draw,minimum width=4.6cm,minimum height=0.9cm] (c) at (0,0) {协调者进程};
\node[draw,inner sep=2pt,fill=white] at (1.75,0.45) {全局事务};
\fill (0,-0.45) circle (2.2pt);
\node[draw,minimum width=1.9cm,minimum height=0.8cm] (a1) at (-1.35,-2.0) {代理者进程 1};
\node[draw,inner sep=2pt,fill=white] at (-1.35,-2.45) {子事务 1};
\node[draw,minimum width=1.9cm,minimum height=0.8cm] (an) at (1.35,-2.0) {代理者进程 $n$};
\node[draw,inner sep=2pt,fill=white] at (1.35,-2.45) {子事务 $n$};
\node at (0,-2.0) {…};
\draw[fill=white] (-1.35,-1.6) circle (2.2pt); \draw[fill=white] (1.35,-1.6) circle (2.2pt);
\draw[-{Stealth[length=4pt]}] (0,-0.45) -- (-1.3,-1.52);
\draw[-{Stealth[length=4pt]}] (0,-0.45) -- (1.3,-1.52);
\fill (-1.6,-3.0) circle (2.2pt); \node[right] at (-1.5,-3.0) {根代理};
\draw[fill=white] (0.3,-3.0) circle (2.2pt); \node[right] at (0.4,-3.0) {子代理};
\end{tikzpicture}
\end{minipage}\hfill
\begin{minipage}[t]{0.6\linewidth}\vspace{0pt}
进程模型：全局事务为每个子事务在相应场地上创建一个代理者进程（子代理），由它执行该场地上的操作；同时启动一个协调者进程（根代理），负责代理者进程之间的通信，控制和协调它们的操作。根代理所在的场地是事务的源场地。只有根代理能创建子代理；所有子事务都成功后，根代理才决定全局提交。\\[2pt]
实例：银行转账，ACCOUNT(Acc\_No, Amount) 中贷方账户 from\_acc 在场地1，借方账户 to\_acc 在场地2，根代理在场地1 \kj{7-23--25, 7-43--44}。
\end{minipage}\\[3pt]
\cir{1} 场地1 的根代理向本地 DTM 发 begin\_transaction，全局事务开始；\cir{2} DTM 向本地 LTM 发 local\_begin，LTM 写日志后执行局部操作：查 from\_acc 的余额，不足则 Abort，否则扣减 transfer\_amount；\cir{3} 根代理用 create 原语创建子代理 agent1；\cir{4} 场地1 的 DTM 通知 agent1 所在场地2 的 DTM，后者向本地 LTM 发 local\_begin；\cir{5} 场地2 的 LTM 写日志；\cir{6} LTM 激活 agent1；\cir{7} agent1 用 receive 接收参数 (to\_acc, transfer\_amount)，执行 \texttt{UPDATE ACCOUNT SET Amount=Amount+transfer\_amount WHERE Acc\_No=to\_acc}；\cir{8} 根代理执行 Commit，用两阶段提交协议提交两个场地上的子事务。
\end{ans}
\end{keep}

\Needspace{14\baselineskip}
\begin{keep}
\begin{prob}[2018 年第三题]
存在如下信息：场地1：货品信息表 Article(Ano, Aname, Price, Sold, Total)，用户信息表 User(Uid, name, Tel, Type)；场地2：购买信息表 Purchase(Pid, Uid, Ano, Money)；场地3：电商账户 B\_account；场地4：用户账号 U\_account。假设用户 A 在某电商的一次购物为一个分布式事务，具体如下：步骤1：客户 A 购买了商品 X，修改该货品 X 的售出数量 Sold；步骤2：写用户 A 的购买信息 Y 到购买信息表 Purchase；步骤3：将用户 A 购买商品 X 的花费从用户账号 U\_account 转账到电商账户 B\_account。请回答：1. 将该分布式事务分解为相应的子事务；2. 应用进程代理模型概述该业务流程的实现过程；3. 如果采用集中式两段提交协议，写出提交过程。
\end{prob}
\begin{ans}[2018 年第三题]
设 U\_account(Uid, Balance)、B\_account(Bid, Balance)，Total 为库存总数，可售数量为 Total$-$Sold。\\
1. 子事务：$T_1$（场地1）：查 X 的单价和库存，修改 X 的 Sold；$T_2$（场地2）：向 Purchase 插入 A 的购买记录；$T_3$（场地3）：电商账户 B\_account 加上花费；$T_4$（场地4）：检查 A 的余额，从 U\_account 扣除花费。\\
2. 根代理在场地1（商品信息所在场地），三个子代理分别在场地4、2、3：\\[1pt]
\begin{minipage}[t]{0.55\linewidth}
\textbf{根代理（场地1）}
\begin{code}
begin_transaction
 input(uid, ano, num, bid);
 SELECT Price,Sold,Total INTO :p,:s,:t
   FROM Article WHERE Ano=:ano;
 IF t-s < num THEN
   Output("\zh{库存不足}"); Abort;
 ELSE
   money := p*num;
   UPDATE Article SET Sold=Sold+:num
     WHERE Ano=:ano;
   Create AGENT1; Send(AGENT1, uid, money);
   Create AGENT2; Send(AGENT2, uid, ano, money);
   Create AGENT3; Send(AGENT3, bid, money);
   Commit;
end
\end{code}
\end{minipage}\hfill
\begin{minipage}[t]{0.43\linewidth}
\textbf{AGENT1（场地4）}
\begin{code}
Receive(ROOT, uid, money);
SELECT Balance INTO :b
  FROM U_account WHERE Uid=:uid;
IF b < money THEN Abort;
UPDATE U_account
  SET Balance=Balance-:money
  WHERE Uid=:uid;
\end{code}
\textbf{AGENT2（场地2）}
\begin{code}
Receive(ROOT, uid, ano, money);
INSERT INTO Purchase(Uid,Ano,Money)
  VALUES(:uid, :ano, :money);
\end{code}
\textbf{AGENT3（场地3）}
\begin{code}
Receive(ROOT, bid, money);
UPDATE B_account
  SET Balance=Balance+:money
  WHERE Bid=:bid;
\end{code}
\end{minipage}\\[2pt]
3. 协调者是场地1 的根代理，参与者是场地4、2、3 的子代理（场地1 本地的子事务由协调者一并准备和提交）。决定阶段：协调者写“开始提交”记录，向三个参与者发“预提交”；参与者检查本地子事务，能提交就写“准备提交”记录并回答“准备提交”，否则写“准备废弃”并回答“准备废弃”（如 AGENT1 余额不足）。执行阶段：全部回答“准备提交”时，协调者写“决定提交”，发“全局提交”，否则写“决定废弃”，发“全局废弃”；参与者写“提交”或“废弃”记录，提交或撤销子事务，回 Ack；协调者收齐 Ack 后写“事务结束”。共 $4\times3=12$ 条消息，流程图见 \S3.2。
\end{ans}
\end{keep}

\Needspace{14\baselineskip}
\begin{keep}
\begin{prob}[2008 年第 3 题]
完成电话银行代缴费业务。设银行账户关系为 ACCOUNT(Acc\_no, Amount1)，Acc\_no 表示银行存款账号，Amount1 表示银行存款余额；通信公司数据库中的客户关系为 CUSTOMER(Phone\_no, Amount2)，Phone\_no 表示客户电话号码，Amount2 表示话费余额。具体的业务过程是：电话用户到达银行，填写拟缴费的电话号码和缴费金额，银行先创建从电话用户账号（设为 111）到通信公司存款账号（设为 222）的转账业务，同时将缴费信息发送给通信公司。通信公司接收到缴费信息后，为该电话号码修改话费余额。假设存放账号 111 的信息表、存放账号 222 信息的表以及通信公司客户关系表均不在一个场地上，请按照进程代理模型给出分布式事务的执行过程。
\end{prob}
\begin{ans}[2008 年第 3 题]
根代理在存放账号 111 的场地（用户所在银行），AGENT1 在存放账号 222 的场地，AGENT2 在通信公司场地。\\[1pt]
\begin{minipage}[t]{0.56\linewidth}
\textbf{根代理（账号 111 所在场地）}
\begin{code}
begin_transaction
 input(from_no, to_no, phone_no, money);
 SELECT Amount1 INTO :temp FROM ACCOUNT
   WHERE Acc_no=:from_no;
 IF temp < money THEN
   Output("\zh{余额不足}"); Abort;
 ELSE
   UPDATE ACCOUNT SET Amount1=Amount1-:money
     WHERE Acc_no=:from_no;
   Create AGENT1; Send(AGENT1, to_no, money);
   Create AGENT2; Send(AGENT2, phone_no, money);
   Commit;
   Output("\zh{缴费完成}");
end
\end{code}
\end{minipage}\hfill
\begin{minipage}[t]{0.42\linewidth}
\textbf{AGENT1（账号 222 所在场地）}
\begin{code}
Receive(ROOT, to_no, money);
UPDATE ACCOUNT
  SET Amount1=Amount1+:money
  WHERE Acc_no=:to_no;
\end{code}
\textbf{AGENT2（通信公司场地）}
\begin{code}
Receive(ROOT, phone_no, money);
UPDATE CUSTOMER
  SET Amount2=Amount2+:money
  WHERE Phone_no=:phone_no;
\end{code}
\end{minipage}\\[2pt]
根代理执行 Commit 时，以两阶段提交协议提交三个场地的子事务：扣款、转入和修改话费余额要么都生效，要么都不生效。
\end{ans}
\end{keep}

\begin{keep}
\begin{prob}[2011 年第 5 题]
完成网上购买飞机票业务，过程是：买票人登录航空公司网站，查询指定日期、指定航班的余票信息。如果有剩余机票，则买票人输入身份信息和机票信息从而生成订单，并进行网上支付。支付成功后，买票过程完毕。现假设航空公司数据、航空公司银行账户数据、买票人银行账户数据分布在三个场地上。银行账户关系为 ACCOUNT(Acc\_no, Amount)，Acc\_no 表示账号，Amount 表示银行存款余额；航空公司数据库中客户关系为 CUSTOMER(ID, Name, Fno)，ID 为乘客的身份证号码，Name 为乘客姓名，Fno 为航班号；FLIGHT(Fno, Fdate, Price, STSold, Capacity)，Fno 为航班号，Fdate 为起飞日期，Price 为航班价格，STSold 为已卖座位数，Capacity 为总座位数。请按照进程代理模型给出分布式事务的执行过程。
\end{prob}
\begin{ans}[2011 年第 5 题]
根代理在航空公司场地，AGENT1 在买票人银行场地，AGENT2 在航空公司银行场地。\\[1pt]
\begin{minipage}[t]{0.58\linewidth}
\textbf{根代理（航空公司场地）}
\begin{code}
begin_transaction
 input(fdate, fno);
 SELECT Capacity-STSold, Price INTO :left,:price
   FROM FLIGHT WHERE Fno=:fno AND Fdate=:fdate;
 IF left <= 0 THEN
   Output("\zh{无余票}"); Abort;
 ELSE
   input(id, name, buyer_acc, air_acc);
   Create AGENT1; Send(AGENT1, buyer_acc, price);
   Create AGENT2; Send(AGENT2, air_acc, price);
   UPDATE FLIGHT SET STSold=STSold+1
     WHERE Fno=:fno AND Fdate=:fdate;
   INSERT INTO CUSTOMER VALUES(:id,:name,:fno);
   Commit;
   Output("\zh{购票成功}");
end
\end{code}
\end{minipage}\hfill
\begin{minipage}[t]{0.4\linewidth}
\textbf{AGENT1（买票人银行）}
\begin{code}
Receive(ROOT, buyer_acc, price);
SELECT Amount INTO :b FROM ACCOUNT
  WHERE Acc_no=:buyer_acc;
IF b < price THEN Abort;
UPDATE ACCOUNT
  SET Amount=Amount-:price
  WHERE Acc_no=:buyer_acc;
\end{code}
\textbf{AGENT2（航空公司银行）}
\begin{code}
Receive(ROOT, air_acc, price);
UPDATE ACCOUNT
  SET Amount=Amount+:price
  WHERE Acc_no=:air_acc;
\end{code}
\end{minipage}\\[2pt]
余票不足时根代理直接 Abort；买票人余额不足时 AGENT1 废弃本地子事务，在两阶段提交中回答“准备废弃”，全局事务废弃，已改的 STSold 和 CUSTOMER 记录一起撤销。
\end{ans}
\end{keep}

% ---------------------------------------------------------------
\Needspace{14\baselineskip}
\subsection{集中式两阶段提交协议 2PC \kj{7-29--34, 7-47}}\label{s:2pc}
\begin{center}
\begin{tikzpicture}[x=1cm,y=0.62cm,font=\footnotesize,
  msg/.style={-{Stealth[length=4pt]},thick},
  lgl/.style={draw,fill=white,inner sep=1.5pt,font=\scriptsize,text width=4.1cm,align=right,anchor=east},
  lgr/.style={draw,fill=white,inner sep=1.5pt,font=\scriptsize,text width=4.1cm,align=left,anchor=west},
  fp/.style={font=\small\bfseries}]
\draw[very thick] (0,0.7) node[above]{协调者 C（根代理）} -- (0,-8.6);
\draw[very thick] (7.4,0.7) node[above]{参与者 P$_i$（子代理），共 $N$ 个} -- (7.4,-8.6);
\node[lgl] at (-0.15,0) {写“开始提交”Begin\_commit};
\draw[msg] (0,-0.4) -- node[above,sloped]{\cir{1} 预提交 Prepare} (7.4,-1.5);
\node[lgr] at (7.55,-1.9) {能提交：写“准备提交”Ready\\不能提交：写“准备废弃”Abort};
\draw[msg] (7.4,-2.5) -- node[above,sloped]{\cir{2} 准备提交 Ready / 准备废弃 Abort} (0,-3.6);
\node[lgl] at (-0.15,-4.0) {全是 Ready：写“决定提交”\\有 Abort 或超时：写“决定废弃”};
\draw[msg] (0,-4.6) -- node[above,sloped]{\cir{3} 全局提交 Commit / 全局废弃 Abort} (7.4,-5.7);
\node[lgr] at (7.55,-6.0) {写“提交”或“废弃”记录\\提交或撤销子事务};
\draw[msg] (7.4,-6.5) -- node[above,sloped]{\cir{4} 确认 Ack} (0,-7.6);
\node[lgl] at (-0.15,-7.9) {收齐 Ack：写“事务结束”};
\draw[decorate,decoration={brace,amplitude=4pt,mirror}] (-4.6,0.3) -- node[left=4pt,align=center]{决\\定\\阶\\段} (-4.6,-4.3);
\draw[decorate,decoration={brace,amplitude=4pt,mirror}] (-4.6,-4.5) -- node[left=4pt,align=center]{执\\行\\阶\\段} (-4.6,-8.3);
\node[fp] at (7.0,-1.05) {\bcir{1}\cross};
\node[fp] at (7.0,-3.35) {\bcir{2}\cross};
\node[fp] at (7.0,-7.1) {\bcir{3}\cross};
\node[fp] at (0.4,-1.7) {\cross\bcir{4}};
\node[fp] at (0.4,-6.0) {\cross\bcir{5}};
\node[font=\scriptsize,align=center] at (3.7,-9.35) {参与者状态：初始 →（发出 R/A 后）准备就绪或废弃 →（发出 Ack 后）提交或废弃\\集中式 2PC 的消息数为 $4N$，耗时 $4T$；\bcir{1}–\bcir{5} 是故障点，处理方法见 \S3.3 的表};
\end{tikzpicture}
\end{center}

\begin{keep}
\begin{prob}[2004 年第 2 题]
画出集中式 2PC 的通信协议结构，编写集中式的 2PC 提交协议的协调者和参与者的执行算法。
\end{prob}
\begin{ans}[2004 年第 2 题]
通信结构见上图：协调者与每个参与者之间一问一答，参与者之间不通信。\\
协调者。决定阶段：写“开始提交”记录，向所有参与者发“预提交”，等待应答。执行阶段：若有一个“准备废弃”或超时没有收到应答，写“决定废弃”，发“全局废弃”；否则写“决定提交”，发“全局提交”。然后等待所有 Ack，收齐后写“事务结束”。\\
参与者。决定阶段：收到“预提交”后，判断本地子事务能否提交。能提交就写“准备提交”记录，回答“准备提交”；否则写“准备废弃”记录，回答“准备废弃”。然后等待协调者的决定。执行阶段：收到“全局提交”就写“提交”记录并提交子事务；收到“全局废弃”就写“废弃”记录并撤销子事务。完成后向协调者发 Ack。\\
两条规则：先写日志，再发消息；参与者回答“准备提交”以后，不能自行废弃，只能等协调者的决定。
\end{ans}
\end{keep}
{\small
\begin{tabularx}{\linewidth}{@{}l X l l@{}}
\toprule
2PC 的实现方法 \kj{7-47--53} & 通信结构 & 消息数 & 耗时\\ \midrule
集中式 & 协调者与参与者一问一答，参与者之间不通信 & $4N$ & $4T$\\
分布式 & 始发场地广播预提交，各场地互相广播应答，各自做决定；只有决定阶段 & $N(N+1)$ & $2T$\\
分层式 & 树形结构，逐层转发，协调者是根 & $4N$ & $4T\sim2NT$\\
线性 & 场地排成链，预提交逐个向后传，提交从链尾逐个传回 & $\le3(N-1)$ & $\le2(N-1)T$\\
\bottomrule
\end{tabularx}}

\Needspace{14\baselineskip}
\subsection{2PC 的故障与处理 \kj{8-42--55}}\label{s:fail}
{\small
\begin{tabularx}{\linewidth}{@{}c >{\raggedright\arraybackslash}p{4.6cm} >{\raggedright\arraybackslash}p{3.6cm} X@{}}
\toprule
故障点 & 故障发生的时机 & 现象 & 处理\\ \midrule
\bcir{1} & 参与者在写 R/A 记录之前 & 协调者等应答超时 & 协调者按废弃处理，视同收到 Abort，全局废弃\\
\bcir{2} & 参与者写了 R/A，还没写 C/A：\par (a) 应答没发出 \par (b) 应答已发出 & (a) 协调者收不齐应答 \par (b) 协调者收不齐 Ack & (a) 参与者恢复后写 A 记录，发 A 应答，全局废弃 \par (b) 协调者重发 C/A 命令，参与者恢复后按命令提交或废弃\\
\bcir{3} & 参与者写了 C/A：\par (a) Ack 没发出 \par (b) Ack 已发出 & (a) 协调者收不齐 Ack \par (b) 正常 & (a) 协调者重发 C/A，参与者按日志 redo 或 undo 后回 Ack \par (b) 不需要处理\\
\bcir{4} & 协调者写了“开始提交”，还没写 C/A：\par (a) P 未发出 \par (b) P 已发出 & (a) 无 \par (b) 参与者已回 R/A，等命令超时 & (a) 协调者重启后重新执行 2PC \par (b) 参与者启动终结协议\\
\bcir{5} & 协调者写了 C/A，还没写结束记录：\par (a) 命令未发出 \par (b) 命令已发出 & (a) 参与者等命令超时 \par (b) 可能漏收 Ack & (a) 参与者启动终结协议 \par (b) 协调者重启后向所有参与者重发 C/A\\ \midrule
通信 & 丢失 P 报文或 R/A 应答 & 协调者等应答超时 & 协调者按废弃处理\\
通信 & 丢失 C/A 命令 & 参与者等命令超时 & 参与者请求协调者重发，或启动终结协议\\
通信 & 丢失 Ack & 协调者等 Ack 超时 & 协调者重发 C/A 命令，要求应答\\
\bottomrule
\end{tabularx}}
\begin{var}[终结协议 \kj{8-52--53}：参与者等不到协调者的命令时使用]
\small 选一个参与者 $P_T$ 作新的协调者，向所有参与者询问状态，然后：有的处于“初始”、有的处于“准备就绪”，发 Abort；全部“准备就绪”，发 Commit；至少一个“提交”，发 Commit；至少一个“废弃”，发 Abort。依据是不会出现“初始”和“提交”并存的情况。
\end{var}
\begin{dlg}
\Q{协调者收不到某个参与者的 Ready 时可以直接废弃，收不到 Ack 时却只能一直重发，为什么？}
\A{收 Ready 时还在决定阶段，没有人做出最终决定，废弃是安全的。收 Ack 时已经在执行阶段，决定已写入日志并发给了其他参与者，有的参与者可能已经提交，不能再改，只能重发命令，直到对方恢复并确认。}
\Q{参与者写了 Ready 以后宕机，恢复后能自己决定提交或废弃吗？}
\A{如果应答还没发出，可以写 Abort（见 \bcir{2}(a)）。应答已经发出时就不能自己决定：它处于不确定状态，必须问协调者，或者问其他参与者（终结协议）。这就是 2PC 的阻塞问题，3PC 在中间加了一个“准备提交”阶段来减少阻塞 \kj{7-55--64}。}
\end{dlg}

\newcommand{\twopc}[1]{%
\begin{tikzpicture}[x=0.62cm,y=0.5cm,font=\scriptsize,
  nd/.style={circle,draw,inner sep=0.6pt,minimum size=11pt},
  ar/.style={-{Stealth[length=3pt]}}]
\node[nd] (c1) at (0,0) {C};
\foreach \i/\y in {1/2,2/1,3/0,4/-1,5/-2}{\node[nd] (a\i) at (3,\y) {P\i}; \draw[ar] (c1)--(a\i);}
\node[nd] (c2) at (6,0) {C};
\foreach \i in {1,...,5}{\draw[ar] (a\i)--(c2);}
\foreach \i/\y in {1/2,2/1,3/0,4/-1,5/-2}{\node[nd] (b\i) at (9,\y) {P\i}; \draw[ar] (c2)--(b\i);}
\node[nd] (c3) at (12,0) {C};
\foreach \i in {1,...,5}{\draw[ar] (b\i)--(c3);}
\node at (1.5,2.6) {预提交}; \node at (4.5,2.6) {准备提交/废弃}; \node at (7.5,2.6) {全局提交/废弃}; \node at (10.5,2.6) {确认 Ack};
\draw[decorate,decoration={brace,amplitude=3pt,mirror}] (0,-2.7) -- node[below=2pt]{决定阶段} (6,-2.7);
\draw[decorate,decoration={brace,amplitude=3pt,mirror}] (6,-2.7) -- node[below=2pt]{执行阶段} (12,-2.7);
#1
\end{tikzpicture}}

\Needspace{14\baselineskip}
\begin{keep}
\begin{prob}[2019 年第 3 题（16 分）]
假设存在如下信息。场地1：电子商品信息表 Article(Ano, Aname, Price, Sold, Total)；场地2：电子附件商品信息表 Article(Ano, Aname, Price, Sold, Total)；场地3：购买信息表 Purchaser(Pid, Uid, Ano, Num, Money)，用户信息表 User(Uid, name, Tel, Type)；场地4：电商账户（B\_account）；场地5：用户账号（U\_account）。假设用户 A 在某电商的一次购物简化为如下一个分布式事务，具体过程：步骤1：客户 A 购买了商品 X（场地1）和商品 Y（场地2），修改商品 X 和 Y 的售出数量 (Sold)；步骤2：写用户 A 的购买信息 XY 到购买信息表 (Purchaser)；步骤3：将用户 A 购买商品 X 和 Y 的花费 Z 从用户账号 (U\_account) 转账到电商账户 (B\_account) 中。请回答：1) 将该分布式事务分解为相应的子事务；2) 如果采用集中式两阶段提交协议，写出该分布式事务的提交过程；3) 在提交过程中，如果转账没有成功，其提交过程如何？请画图说明故障点，并给出处理策略；4) 如果在发 commit 命令后场地3 出故障，其提交过程又如何？请画图说明故障点，并给出处理策略。
\end{prob}
\begin{ans}[1) 分解子事务 \quad 2) 集中式 2PC 的提交过程]
设事务由电商平台的源场地 $S_0$ 发起，根代理（协调者）在 $S_0$。全局事务分解为 5 个子事务：
$T_1$（场地1）：检查 X 的库存，修改 X 的 Sold；
$T_2$（场地2）：检查 Y 的库存，修改 Y 的 Sold；
$T_3$（场地3）：向 Purchaser 插入 A 购买 X、Y 的记录；
$T_4$（场地4）：B\_account 余额加 $Z$；
$T_5$（场地5）：检查 A 的余额，U\_account 余额减 $Z$。\\[2pt]
提交过程：协调者 C 是 $S_0$ 的根代理，参与者 P1–P5 是场地1–5 的子代理（下方左图）。
决定阶段：C 写“开始提交”，向 P1–P5 发“预提交”。各 $P_i$ 检查本地子事务（P1、P2 库存够，P5 余额够，P3、P4 写入成功），能提交就写“准备提交”并回答“准备提交”，否则写“准备废弃”并回答“准备废弃”。
执行阶段：C 收到 5 个“准备提交”，写“决定提交”，发“全局提交”；各 $P_i$ 写“提交”记录，提交子事务，回 Ack；C 收齐 5 个 Ack，写“事务结束”。共 $4\times5=20$ 条消息。\\[3pt]
\centerline{\twopc{}\qquad
\twopc{\draw[ultra thick] (a5)--(c2); \node[font=\small\bfseries] at (4.6,-1.6) {\cross}; \node[font=\scriptsize,fill=white,inner sep=1pt] at (4.6,-2.15) {P5 答“准备废弃”}; \node[font=\scriptsize] at (7.5,-1.6) {全局废弃};}}
\centerline{\footnotesize 左：正常提交 \hspace{5.2cm} 右：3) 转账失败，P5 余额不足}
\end{ans}
\end{keep}
\begin{ans}[3) 转账没有成功]
故障点：决定阶段，P5（U\_account）发现余额小于 $Z$（或者 P4、P5 执行转账失败），本地子事务废弃，见上方右图。\\
提交过程：P5 写“准备废弃”记录，回答“准备废弃”（P5 宕机没有应答也一样，C 超时后按废弃处理）。C 收到“准备废弃”，写“决定废弃”，向 P1–P5 发“全局废弃”。各参与者写“废弃”记录，撤销本地子事务（P1、P2 恢复 Sold，P3 删除购买记录，P4 恢复余额），回 Ack。C 收齐 Ack 后写“事务结束”。全局事务废弃，所有场地回到购物前的状态。
\end{ans}
\begin{ans}[4) 发出 commit 命令后场地3 出故障]
\begin{minipage}[t]{0.5\linewidth}\vspace{0pt}
故障点：C 已写“决定提交”并向 P1–P5 发出“全局提交”，P3 宕机（右图）。P1、P2、P4、P5 照常提交并回 Ack，C 收不到 P3 的 Ack。\\
处理方法取决于 P3 宕机时日志写到了哪里：
\end{minipage}\hfill
\begin{minipage}[t]{0.47\linewidth}\vspace{0pt}\centering
\twopc{\node[font=\small\bfseries] at (9.75,0.55) {\cross}; \draw[dashed,thick] (b3)--(c3); \node[font=\scriptsize,fill=white,inner sep=1pt] at (10.5,-0.75) {P3 宕机，无 Ack};}
\end{minipage}\\[2pt]
\cir{1} P3 还没写“提交”记录，日志里只有“准备提交”：C 超时后反复重发“全局提交”。P3 恢复后发现子事务处于“准备就绪”的不确定状态，不能自己决定，要向 C 询问，或者向其他参与者询问（终结协议，此时已有参与者提交，所以提交）。得到 Commit 后写“提交”记录，按日志 redo 插入的购买记录并提交，回 Ack。\\
\cir{2} P3 已写“提交”记录，还没回 Ack：恢复时把它作为已提交事务 redo；收到 C 重发的命令后回 Ack。\\
\cir{3} P3 已回 Ack：不需要处理。\\
全局事务仍然提交，其他场地不受影响。C 收齐包括 P3 在内的所有 Ack 以后才写“事务结束”。
\end{ans}
\begin{dlg}
\Q{如果场地3 正好是协调者所在场地呢？}
\A{那是故障点 \bcir{5}。协调者写了“决定提交”并发出命令后宕机，恢复后读日志，看到“决定提交”而没有“事务结束”，就向所有参与者重发“全局提交”，收齐 Ack 后写结束记录。如果宕机时命令还没发出，参与者等待超时后启动终结协议。}
\end{dlg}

\Needspace{14\baselineskip}
\begin{keep}
\begin{prob}[2020 年第 4 题（9 分）]
某公司根据业务需求设置了 4 个部门，分别为总部、商品研发部、商品生产部和商品销售部，分别在沈阳（场地0）、北京（场地1）、上海（场地2）和广州（场地3）四个城市。原题附图：场地0 总部存“雇员”，场地1 商品研发部存“研发型商品”，场地2 商品生产部存“生产型商品”，场地3 商品销售部存“销售明细”。假设存在如下全局关系模式：雇员（雇员号，姓名，性别，部门）；商品（商品号，商品名，品牌，商品类型，商品负责人），商品类型分为研发型和生产型，对应的负责人为研发负责人和生产负责人；销售明细（商品号，销售数量，销售人员）。假设总部发布一个全局事务 T1：\\
\hspace*{1em}\texttt{Begin transaction T1\ \ Select \zh{商品类型}, count(*) from \zh{商品} group by \zh{商品类型}\ \ End}\\
请回答：1) 执行该全局事务时，请说明：该全局事务分为几个子事务；定义协调者和参与者，并说明它们完成的功能；2) 若采用集中式两阶段提交协议提交，简述提交过程；3) 在提交过程中，如果丢失了 commit 命令，如何处理？
\end{prob}
\begin{ans}[2020 年第 4 题]
1) 商品按商品类型水平分片，研发型存场地1，生产型存场地2，所以 T1 分为 2 个子事务：
$T_{11}$（场地1）：\texttt{SELECT '\zh{研发型}', count(*) FROM \zh{研发型商品}}；$T_{12}$（场地2）：\texttt{SELECT '\zh{生产型}', count(*) FROM \zh{生产型商品}}。\\
协调者：总部场地0 的根代理（事务的源场地）。它分解全局事务，创建子代理并发送子查询，收集两个计数，合并成分组结果返回给用户，然后发起 2PC，决定全局提交或废弃。\\
参与者：场地1、场地2 的子代理。它们执行本地子事务，返回计数；在 2PC 中报告“准备提交”或“准备废弃”，按协调者的命令提交或废弃，并回 Ack。\\
2) 协调者写“开始提交”，向场地1、2 发“预提交”；两个参与者写“准备提交”并应答；协调者写“决定提交”，发“全局提交”；参与者写“提交”记录、提交、回 Ack；协调者写“事务结束”。共 $4N=8$ 条消息，耗时 $4T$。（T1 只读，实际系统可以对只读参与者做优化，答题时按标准流程写。）\\
3) 参与者回答“准备提交”后收不到命令，处于等待状态。超时后，它请求协调者重发命令，或者启动终结协议：询问另一个参与者的状态，对方已提交就提交，双方都是“准备就绪”也提交。协调者收不到该参与者的 Ack，超时后重发 Commit，直到收到 Ack 再写“事务结束”。
\end{ans}
\end{keep}

\clearpage
\section{并发控制：冲突、局部与全局可串行化、2PL}

\subsection{判定与构造的方法 \kj{9-10--18}}\label{s:cc}
\begin{flow}[可串行化判定与构造的步骤]
\begin{enumerate}
\item 把每个事务写成 $R_i(x)$、$W_i(x)$ 序列，并标出场地。“修改 $x$”写成 $R_i(x)W_i(x)$；“插入”是写；“查询、统计”是读。
\item 找冲突对：同一场地、同一数据项、不同事务、至少一个是写，即 R–W、W–R、W–W；R–R 不冲突。
\item 判断局部可串行：在每个场地的局部历程 $H(S_k)$ 里，冲突对 $O_i<O_j$ 给出一条边 $T_i\to T_j$。没有环就可串行，按拓扑序写出 $SR(H)$。
\item 判断全局可串行：每个局部历程都可串行，并且存在一个与所有场地的边都一致的总序；等价于把各场地的边画在一张图上没有环 \kj{9-17}。
\item 构造“局部可串行、全局不可串行”：找两个在两个场地上都有冲突的事务 $T_a,T_b$，让场地 1 上 $T_a$ 在前、场地 2 上 $T_b$ 在前。每个场地内部写成串行的，局部可串行就很明显。
\end{enumerate}
\end{flow}
\begin{dlg}
\Q{为什么由冲突对决定顺序？}
\A{两个不冲突的相邻操作交换位置，不影响结果；冲突操作交换以后，读到的值或最后写入的值会变。所以只要每对冲突操作的先后与某个串行历程一致，两个历程就等价（定理 1 的引理 \kj{9-15}）。}
\Q{每个场地都可串行了，全局为什么还可能不可串行？}
\A{各场地可能对应不同的顺序。比如场地 1 等价于 $T_1T_3$，场地 2 等价于 $T_3T_1$。每个场地单看都没问题，但没有一个全局顺序能同时满足两者，$T_3$ 看到的是做了一半的 $T_1$。}
\Q{构造反例需要什么条件？}
\A{合并后的图要有环，至少需要两个事务在两个不同场地上都有冲突。如果冲突关系是星形的，例如 $T_1$ 和 $T_2$ 只在场地 1 冲突、$T_1$ 和 $T_3$ 只在场地 2 冲突，合并后的图不会有环，怎么排都构造不出来。另外要注意时间上能实现：同一事务里先发生的操作，不能因为在另一个场地的排法而变成后发生。}
\Q{数据项按元组算还是按属性算？}
\A{按封锁粒度定，一般写“以元组（记录）为数据项”。按元组算时，“改 X 的 Sold”和“读 X 的 Total”冲突，因为是同一条记录 X；按属性算则不冲突。题目要求构造全局不可串行的历程时，通常要按元组算。}
\end{dlg}

\subsection{2023 年第五题}
\begin{keep}
\begin{prob}[原题]
设数据项 $x$，$y$ 存放在 S1 场地，$u$，$v$ 存放在 S2 场地。有分布式事务 $T_1$ 和 $T_2$，判断下面的每个执行是否是局部可串行的？并分别说明理由。在 S1 场地执行历程 $H_1$：$R_1(x)R_1(y)W_1(x)R_2(x)W_1(y)W_2(x)$；在 S2 场地执行历程 $H_2$：$R_1(u)R_1(v)R_2(u)W_1(u)W_1(v)W_2(u)$。
\end{prob}
\begin{ans}[2023 年第五题]
$H_1$：$x$ 上的冲突对是 $R_1(x)<W_2(x)$、\allowbreak $W_1(x)<R_2(x)$、\allowbreak $W_1(x)<W_2(x)$，都是 $T_1\to T_2$；$y$ 只被 $T_1$ 访问。没有环，$H_1$ 是可串行的，等价于串行历程 $T_1T_2$：\\
\hspace*{2em}$R_1(x)R_1(y)W_1(x)W_1(y)R_2(x)W_2(x)$。\\
$H_2$：$u$ 上的冲突对是 $R_1(u)<W_2(u)$（$T_1\to T_2$）、\allowbreak $R_2(u)<W_1(u)$（$T_2\to T_1$）、\allowbreak $W_1(u)<W_2(u)$（$T_1\to T_2$）。出现 $T_1\to T_2\to T_1$ 的环，$H_2$ 不可串行：它不等价于 $T_1T_2$（$R_2(u)<W_1(u)$ 的次序不同），也不等价于 $T_2T_1$（$R_1(u)<W_2(u)$ 的次序不同）。因此这组分布式事务的执行也不是全局可串行的。
\end{ans}
\end{keep}

\Needspace{14\baselineskip}
\subsection{2019 年第 4 题（8 分）}
\begin{keep}
\begin{prob}[原题]
针对题 3 信息，假设存在如下三个并发事务：事务 $T_1$ 为客户 A 购买了商品 X（场地1）和商品 Y（场地2），并修改商品 X 和 Y 的售出数量 (Sold)；事务 $T_2$ 为仓储管理者 B 入库商品 X（场地1），并修改商品 X 的数量 (Total)；事务 $T_3$ 为电商经理 C 查询商品 X（场地1）和商品 Y（场地2）的库存数量 (Total)。请回答：1) 几个场地上存在冲突操作？若存在冲突操作，请写出冲突对；2) 针对上述三个分布式事务，请构建一全局历程 H，同时写出：a) 局部可串行化、全局也可串行化的局部历程；b) 局部可串行化、全局不可串行化的局部历程。
\end{prob}
\begin{ans}[2019 年第 4 题]
读写序列（以商品元组为数据项，“修改”是先读后写）：\\
\hspace*{1em}$T_1$：$R_1(X)W_1(X)$（场地1），$R_1(Y)W_1(Y)$（场地2）；\quad $T_2$：$R_2(X)W_2(X)$（场地1）；\quad $T_3$：$R_3(X)$（场地1），$R_3(Y)$（场地2）。\\
1) 场地1 和场地2 上都有冲突操作。\\
\hspace*{1em}场地1：\quad $T_1$–$T_2$：$(R_1(X),W_2(X))$、$(W_1(X),R_2(X))$、$(W_1(X),W_2(X))$；\\
\hspace*{4.4em}$T_1$–$T_3$：$(W_1(X),R_3(X))$；\quad $T_2$–$T_3$：$(W_2(X),R_3(X))$。\\
\hspace*{1em}场地2：\quad $T_1$–$T_3$：$(W_1(Y),R_3(Y))$。$R_1(Y)$ 与 $R_3(Y)$ 都是读，不冲突。\\[2pt]
2a) 局部可串行、全局也可串行：\\
\hspace*{1em}$H(S_1)=R_1(X)W_1(X)R_2(X)W_2(X)R_3(X)$，$T_1<T_2<T_3$；\quad
$H(S_2)=R_1(Y)W_1(Y)R_3(Y)$，$T_1<T_3$。\\
\hspace*{1em}全局历程 $H=R_1(X)W_1(X)R_1(Y)W_1(Y)R_2(X)W_2(X)R_3(X)R_3(Y)$。总序 $T_1<T_2<T_3$ 与两个场地一致，全局可串行。\\[2pt]
2b) 局部可串行、全局不可串行：\\
\hspace*{1em}$H(S_1)=R_1(X)W_1(X)R_2(X)W_2(X)R_3(X)$，串行，$T_1<T_2<T_3$；\quad
$H(S_2)=R_3(Y)R_1(Y)W_1(Y)$，串行，$T_3<T_1$。\\
\hspace*{1em}全局历程 $H=R_3(Y)\,R_1(X)W_1(X)\,R_1(Y)W_1(Y)\,R_2(X)W_2(X)\,R_3(X)$。\\
\hspace*{1em}两个局部历程都可串行，但场地1 要求 $T_1<T_3$，场地2 要求 $T_3<T_1$，不存在一致的总序，所以全局不可串行。C 读到的 X 是 A 购买之后的值，Y 却是购买之前的值，两者不一致。
\end{ans}
\end{keep}

\subsection{2020 年第 5 题（8 分）}
\begin{prob}[原题]
接题 4，如果存在另两个全局事务 $T_2$ 和 $T_3$：\\
\hspace*{1em}\texttt{Begin transaction T2\ \ Insert into \zh{商品} values("010","R010"," ","\zh{研发型}","\zh{王一}")\ \ End}\\
\hspace*{1em}\texttt{Begin transaction T3\ \ Insert into \zh{商品} values("110","B110"," ","\zh{生产型}","\zh{李二}")\ \ End}\\
假设 $T_1$、$T_2$、$T_3$ 并发执行，请回答：1) 指出存在冲突操作的场地，并写出存在的冲突对；2) 构建一个局部可串行化，而全局不可串行化的调度历程，并说明为什么？
\end{prob}
\begin{dlg}
\Q{按字面写读写序列是什么样？}
\A{记 $A$ 为场地1 的研发型商品片段，$B$ 为场地2 的生产型商品片段。统计要读整个片段，所以以片段为数据项。于是 \allowbreak $T_1$：$R_1(A)$、\allowbreak $R_1(B)$；\allowbreak $T_2$：$W_2(A)$；\allowbreak $T_3$：$W_3(B)$。}
\Q{这样能构造出“局部可串行、全局不可串行”吗？}
\A{只看字面不能。$T_1$ 和 $T_2$ 只在场地1 冲突，$T_1$ 和 $T_3$ 只在场地2 冲突，$T_2$ 和 $T_3$ 不冲突。冲突关系是星形的，合并后最多两条边，不会有环。要构造反例，需要补上插入的语义：插入时要检查商品号（主码）全局唯一，所以 $T_2$ 要到场地2 读 $B$ 查重（$R_2(B)$），$T_3$ 要到场地1 读 $A$ 查重（$R_3(A)$）。这样 $T_2$、$T_3$ 在两个场地都冲突，就可能有环。答卷里要写出这个假设。}
\end{dlg}
\begin{ans}[2020 年第 5 题]
1) 冲突场地与冲突对（$A$、$B$ 的含义同上）：场地1：$(R_1(A),W_2(A))$；场地2：$(R_1(B),W_3(B))$。场地0、场地3 没有冲突。\\
\hspace*{1em}如果考虑插入时的主码唯一性检查（$T_2$ 读 $B$，$T_3$ 读 $A$），还有：场地1 $(R_3(A),W_2(A))$，场地2 $(R_2(B),W_3(B))$。\\
2) 构造（在主码检查的假设下；按字面，冲突关系是星形的，不会成环，理由见上面的问答）：\\
\hspace*{1em}$T_1=\{R_1(A),R_1(B)\}$，$T_2=\{R_2(B),W_2(A)\}$，$T_3=\{R_3(A),W_3(B)\}$，$T_2$、$T_3$ 都是先查重、后插入。\\
\hspace*{1em}$H(S_1)=R_3(A)\,R_1(A)\,W_2(A)$：冲突 $R_3(A)<W_2(A)$、$R_1(A)<W_2(A)$，等价于串行 $T_3T_1T_2$，可串行；\\
\hspace*{1em}$H(S_2)=R_2(B)\,R_1(B)\,W_3(B)$：冲突 $R_2(B)<W_3(B)$、$R_1(B)<W_3(B)$，等价于串行 $T_2T_1T_3$，可串行。\\
\hspace*{1em}全局的时间顺序可以是 $R_2(B)\,R_3(A)\,R_1(A)\,R_1(B)\,W_2(A)\,W_3(B)$，每个事务都先读后写，时间上可以实现。\\
\hspace*{1em}全局不可串行的原因：场地1 要求 $T_3<T_2$，场地2 要求 $T_2<T_3$，两者矛盾，不存在满足所有局部先后关系的总序。两个事务都在对方插入之前查重，如果两条记录的主码相同，就会插入重复的主码。
\end{ans}

\subsection{两阶段封锁协议（2023 年第六题）\kj{9-24--33}}\label{s:2pl}
\begin{keep}
\begin{prob}[原题]
叙述两阶段封锁协议 2PC 的基本思想，分别说明使用 2PC 协议进行集中式和分布式并发控制实现方法。
\end{prob}
\begin{ans}[2023 年第六题]
基本思想：把事务的加锁和解锁分成两个阶段。加锁阶段，事务读写数据前必须先加相应的锁（读锁共享，写锁排斥），申请不到相容的锁就等待；解锁阶段，事务一旦释放了某个锁，就不能再申请新的锁。所有事务都遵守两阶段封锁协议时，并发调度是可串行化的。严格的两阶段封锁协议把锁一直保持到 commit 或 abort 之后再释放，可以避免级联废弃。\\[2pt]
集中式实现（整个系统只有一个 2PL 调度器，在中心场地）\kj{9-31}：\cir{1} 协调者 TM 向中心场地的锁管理器 LM 发加锁请求；\cir{2} LM 判断相容性，相容就允许加锁，否则让事务等待；\cir{3} TM 把操作发给各数据场地的数据处理器 DP 执行；\cir{4} DP 返回操作结束；\cir{5} TM 通知 LM 释放锁。优点是实现简单、便于管理；缺点是中心场地负载重，存在单点失效。\\[2pt]
分布式实现（每个场地有自己的 2PL 调度器）\kj{9-32}：\cir{1} 协调者 TM 把加锁请求发给数据所在场地的 LM；\cir{2} 各 LM 加锁成功后，直接把操作交给本场地的 DP 执行；\cir{3} DP 执行完后通知协调者 TM 操作结束；\cir{4} TM 通知各参与场地的 LM 释放锁。优点是没有中心场地的瓶颈和单点失效；缺点是锁管理复杂，可能出现分布式死锁。\\[2pt]
有副本时的封锁方法 \kj{9-33}：主副本法，锁只加在主副本上；主场地法，锁都加在主场地上；读写全法，读锁封锁一个副本，写锁封锁全部副本；多数副本法，读锁和写锁都要封锁一半以上的副本。
\end{ans}
\end{keep}
\begin{pit}
题目写的是“两阶段封锁协议 2PC”，内容是封锁（Locking），一般缩写为 2PL；两阶段提交协议（Commit）才是 2PC。按封锁协议作答，不要写成预提交、全局提交的流程。2023 年参考答案在这里写错了。
\end{pit}

% ================= 第5部分 恢复 =================
\clearpage
\section{故障恢复：日志、检查点、undo 与 redo}\label{s:rec}
\begin{flow}[由日志求 undo 表和 redo 表的步骤]
\begin{enumerate}
\item 记号：$B_i$ 开始，$D^j_i$ 是 $T_i$ 的第 $j$ 个数据操作，$C_i$ 提交，$A_i$ 废弃，$K$ 是检查点 \kj{8-21}。
\item 画并发执行图：每个事务一行，横轴按日志顺序排列，用竖线标出 $K$ 和故障点。
\item 初始化：undo 表 = 检查点时刻的活动事务（在 $K$ 之前开始、在 $K$ 之前没有 C 或 A 的事务），redo 表为空。
\item 从 $K$ 向后扫描到日志末尾：遇到 $B_i$，把 $T_i$ 加入 undo 表；遇到 $C_i$ 或 $A_i$，把 $T_i$ 从 undo 表移到 redo 表（课件 8-39 和课件末尾速查表都按这个口径）。
\item 执行：从日志末尾反向扫描，反做 undo 表中的事务；再从 $K$ 正向扫描，重做 redo 表中的事务。$K$ 之前已经结束的事务不用处理。
\end{enumerate}
\end{flow}
\begin{dlg}
\Q{检查点之前就提交了的事务为什么不用管？}
\A{检查点要做四件事：\cir{1} 把日志缓冲区写到外存日志；\cir{2} 在外存日志里登记检查点记录；\cir{3} 把数据库缓冲区写到外存数据库；\cir{4} 把检查点记录在日志中的地址写入重启动文件 \kj{8-18}。第 \cir{3} 步完成后，检查点之前提交的事务的修改都已写入磁盘，所以恢复只需从最近的检查点开始。}
\Q{有 $A_i$ 记录的事务，为什么放进 redo 表而不是 undo 表？}
\A{$A_i$ 说明故障前这个事务已经做完了回滚，回滚时写的数据也记在日志里，所以它和提交的事务一样已经结束。恢复时按日志把它重做一遍，先做的修改和后做的撤销都重放，数据库就回到故障前的状态。课件 8-39、课件末尾的速查表和 2023 年参考答案都这样处理。课件 8-38 那一页只把它从 undo 表删掉、不放进 redo 表，两页不一致，答题按 8-39 写。}
\Q{undo 和 redo 分别保证什么？}
\A{undo 把未提交事务已写入外存的数据恢复成前像，保证原子性；redo 把已提交但还没写入外存的数据恢复成后像，保证持久性。两者都是幂等的，做几次结果都一样 \kj{8-22--23}。}
\end{dlg}

\begin{keep}
\begin{prob}[2023 年第二题]
下面是当一个数据库系统出现故障时，日志文件中的信息：\quad
$B_1\ \ B_2\ \ D^1_1\ \ D^1_2\ \ C_1\ \ K\ \ B_3\ \ D^1_3\ \ C_3\ \ D^2_2\ \ B_4\ \ A_5\ \ D^1_4$\ \ Crash\\
根据上述 log 信息，完成下面的处理：1. 画出对应的事务并发执行图；2. 说明检查点的作用和检查点时刻数据库需要完成的主要操作；3. 确定出反做 (undo) 和重做 (redo) 事务集（写出详细过程）；4. 叙述 undo 和 redo 思想，详细写出其基于日志的数据处理过程。
\end{prob}
\begin{ans}[2023 年第二题]
说明：日志中的废弃记录印作 $A_5$，但 $T_5$ 没有开始记录。这道题和课件 8-38、8-39 的例题相同，那里是 $A_2$。下面按 $A_2$ 作答，第 3 问末尾附按 $A_5$ 理解时的结果。\\[2pt]
1. 并发执行图\\[1pt]
{\small\setlength{\tabcolsep}{3.2pt}
\begin{tabular}{@{}c|cccccc|ccccccc|c@{}}
 & 1 & 2 & 3 & 4 & 5 & & 7 & 8 & 9 & 10 & 11 & 12 & 13 & \\ \hline
$T_1$ & $B_1$ & & $D^1_1$ & & $C_1$ & \multirow{4}{*}{$K$} & & & & & & & & \multirow{4}{*}{\rotatebox{90}{\scriptsize Crash}}\\
$T_2$ & & $B_2$ & & $D^1_2$ & & & & & & $D^2_2$ & & $A_2$ & & \\
$T_3$ & & & & & & & $B_3$ & $D^1_3$ & $C_3$ & & & & & \\
$T_4$ & & & & & & & & & & & $B_4$ & & $D^1_4$ & \\
\end{tabular}}\\[3pt]
2. 检查点是在日志中周期性设置的操作标志，作用是减少故障恢复的工作量，恢复时只需从最近的检查点开始。在检查点上要完成：\cir{1} 把日志缓冲区的内容写入外存日志；\cir{2} 在外存日志中登记一个检查点记录；\cir{3} 把数据库缓冲区的内容写入外存数据库；\cir{4} 把外存日志中检查点的地址写入重启动文件。检查点记录里还登记当时的活动事务清单。\\[2pt]
3. \cir{1} 从重启动文件找到最后一个检查点 $K$。$K$ 时刻 $T_1$ 已提交，活动事务只有 $T_2$，初始化 undo 表 = \{$T_2$\}，redo 表 = \{\}。\cir{2} 从 $K$ 正向扫描：\\
\hspace*{1em}{\small\begin{tabular}{@{}lll@{}}
遇到 $B_3$：& $T_3$ 开始，加入 undo 表 & undo = \{$T_2,T_3$\}，redo = \{\}\\
遇到 $C_3$：& $T_3$ 提交，移入 redo 表 & undo = \{$T_2$\}，redo = \{$T_3$\}\\
遇到 $B_4$：& $T_4$ 开始，加入 undo 表 & undo = \{$T_2,T_4$\}，redo = \{$T_3$\}\\
遇到 $A_2$：& $T_2$ 废弃，移入 redo 表 & undo = \{$T_4$\}，redo = \{$T_3,T_2$\}\\
\end{tabular}}\\
\hspace*{1em}\cir{3} 扫描到日志末尾。结论：反做 $T_4$，重做 $T_3$、$T_2$；$T_1$ 在检查点之前已提交，不处理。\\
\hspace*{1em}如果确实是 $A_5$，只能认为 $T_5$ 在这段日志之前就已开始，检查点时处于活动状态：undo 表初值为 \{$T_2,T_5$\}，最后 undo = \{$T_2,T_4$\}，redo = \{$T_3,T_5$\}。\\[2pt]
4. undo（反做、回滚）把数据项恢复到修改之前的值（前像），用于未提交、但修改可能已写入外存的事务，保证原子性。redo（重做、前滚）把数据项恢复到修改之后的值（后像），用于已提交、但修改可能还没写入外存的事务，保证持久性。
基于日志的处理过程（no\_fix/no\_flush 策略）：\cir{1} 局部恢复管理器 LRM 从重启动文件取得最近检查点的地址，建立重做表和反做表，重做表初值为空，反做表初值为检查点上的活动事务；\cir{2} 从检查点向日志末尾扫描，遇到新开始的事务放入反做表，遇到提交或废弃的事务从反做表移到重做表；\cir{3} 反向扫描日志，反做反做表中的事务，直到这些事务的 B 记录；\cir{4} 从检查点正向扫描，重做重做表中的事务，直到相应的 C 记录。
\end{ans}
\end{keep}

{\small
\begin{tabularx}{\linewidth}{@{}l X l@{}}
\toprule
缓冲区更新策略 \kj{8-32--37} & 特点 & 恢复时需要\\ \midrule
fix / flush & 执行中页面固定，未提交的修改不写盘；结束时强制写盘 & 都不需要 \\
fix / no\_flush & 未提交的修改不写盘；已提交的修改不一定写盘 & 只要 redo \\
no\_fix / flush & 未提交的修改可能写盘；已提交的修改一定写盘 & 只要 undo \\
no\_fix / no\_flush & 两种情况都可能出现 & undo 和 redo，即上面的流程 \\
\bottomrule
\end{tabularx}}

\clearpage
\section{数据复制与一致性协议}

\subsection{主从同步复制（2020 年第 3 题）\kj{10-11--28}}\label{s:rep}
\begin{prob}[2020 年第 3 题（8 分）原题]
假设单主站点的情况下，有三个节点 A、B、C，其中 A 保存着 $x$ 和 $y$ 的主副本，B 和 C 保存着它们的副本，考虑下面事务：\quad $T_1$：Read($x$)\quad Write($x$)\quad Commit\\
假设 $T_1$ 在 B 上提交，请回答：1) 结合同步复制协议的执行方式；2) 解释如何保证数据一致性。
\end{prob}
\begin{dlg}
\Q{同步复制和异步复制的区别在哪里？}
\A{看事务提交前是否等所有副本都更新完。同步复制在同一个事务里把更新传到全部副本，提交时所有副本的值相同，也就是读一写全（ROWA），能保证强一致；缺点是慢，受最慢的节点拖累，某个副本坏了事务就完成不了。异步复制先在主副本上更新并提交，之后再刷新其他副本，响应快、可用性高，但可能读到旧数据，还可能出现“事务倒转” \kj{10-11--12}。}
\Q{$T_1$ 在从节点 B 上提交，B 能直接写自己的副本吗？}
\A{不能。单主节点结构下，封锁和更新都先作用在主副本上。B 上的事务管理器 TM 负责协调：读写请求先交给主节点 A 加锁，由 A 确定冲突更新的先后，再由 TM 把写操作发到所有从副本。这是完全复制透明的协议，应用只需把事务交给本地 TM \kj{10-27--28}。}
\end{dlg}
\begin{ans}[2020 年第 3 题]
1) 执行方式（主从结构加同步复制，B 的 TM 作协调者，完全复制透明）：\\
\cir{1} Read($x$)：B 的 TM 把读请求发给主节点 A；A 在主副本 $x_A$ 上加读锁，通知 TM；TM 把 Read($x$) 转给任意一个存有 $x$ 副本的节点执行，一般就在本地 B 读 $x_B$。\\
\cir{2} Write($x$)：TM 把写请求发给 A；A 给 $x_A$ 加写锁，在主副本上执行 Write($x$)，通知 TM 已获得写锁；TM 再把 Write($x$) 发给存有 $x$ 从副本的 B、C，它们在本地副本上执行写操作。冲突的更新在各从节点上的执行顺序与主节点相同。\\
\cir{3} Commit：同步复制要求所有副本一起提交，所以由 B 的 TM 作协调者，A、B、C 作参与者执行 2PC：预提交、准备提交、全局提交、Ack。提交完成后 A 释放 $x$ 上的锁。\\
如果采用有限复制透明，即应用知道主节点的位置：$T_1$ 的读写直接交给 A 执行，A 加锁、写主副本，再由主节点的 TM 把 Write($x$) 发到 B、C，A 担任 2PC 的协调者 \kj{10-25}。\\[2pt]
2) 保证一致性的方法：\cir{1} 对 $x$ 的锁都加在主副本上，冲突操作在 A 上排好先后，各从副本按同样的顺序执行；\cir{2} 写操作在同一个事务内传到 A、B、C 的全部副本（ROWA），读任何一个副本都能读到最新值；\cir{3} 用 2PC 原子提交，所有副本要么都提交新值，要么都不提交，不会只有部分副本更新；\cir{4} 锁保持到提交之后才释放（严格 2PL），其他事务读不到中间状态。所以提交后各副本的值相同，达到强一致性。
\end{ans}
{\small
\begin{tabularx}{\linewidth}{@{}l X X@{}}
\toprule
 & 主从式复制 & 对等式复制\\ \midrule
更新地点 & 只能在主副本上更新，再传播到从副本 & 任一副本都可以更新，再传播到其他副本 \\
优点 & 实现简单，主副本上总是最新值，容易维护一致性 & 没有单点瓶颈，负载均衡，可用性高 \\
缺点 & 主节点负载重，是瓶颈；自治性低 & 同一数据可能在多处同时更新，要做冲突消解，开销大 \\
\bottomrule
\end{tabularx}}

\Needspace{14\baselineskip}
\subsection{Paxos、反熵、NWR 的比较（2018 年第五题）\kj{10-38--54}}\label{s:paxos}
\begin{keep}
\begin{prob}[原题]
对比分析维护多副本一致性协议 Paxos、反熵、NWR 的异同，请写出：各自的优势和不足？分别适合于什么场景？
\end{prob}
\begin{ans}[2018 年第五题]
相同点：三者都通过副本之间传递消息，让多个副本的数据达到一致，都允许部分节点失效。不同点在于一致性的强弱（强一致、最终一致、可配置）和协调方式（多数派投票、节点两两随机交换、读写法定数）。\\[2pt]
{\small
\begin{tabularx}{\linewidth}{@{}l X X X@{}}
\toprule
 & Paxos & 反熵（Gossip 的一种） & NWR \\ \midrule
思想 & 通过消息传递从多个冲突的提议中选出一个。两个阶段：Prepare($n$) 得到多数 acceptor 的承诺，Accept($n,v$) 得到多数批准。角色有 proposer、acceptor、learner & 节点周期性地随机选择其他节点，交换 (key, value, version) 摘要，以 push、pull 或 push-pull 方式按版本号更新 & 共 $N$ 个副本，写操作至少写成功 $W$ 个，读操作至少读 $R$ 个。$W+R>N$ 时读集合和写集合必定相交，至少能读到一个最新版本 \\
一致性 & 强一致，由多数派决定 & 最终一致 & 可配置：$W+R>N$ 时强一致，否则较弱 \\
优势 & 强一致；能容忍少数节点故障和消息丢失、重复（非拜占庭模型） & 去中心化，节点可以随时加入或退出，容错性好，信息按指数速度传播 & 可以按需求设定读写性能和一致性级别；吞吐量高，容易扩展 \\
不足 & 算法复杂，难以实现；需要多轮消息，开销大；不能处理恶意节点 & 有消息延迟和冗余，不能保证实时一致 & 参数配置复杂；$W+R\le N$ 时可能读到旧值；$W=N$ 时写操作的可用性差 \\
场景 & 对一致性要求高的系统，如分布式数据库、选主、锁服务（Chubby、Spanner） & 实时性要求低、最终一致即可的场景，如成员管理、分布式存储和备份、Dynamo 和 Cassandra 的副本同步 & 需要在读写性能和一致性之间折中的键值存储，常取 $N=3$、$W=R=2$ \\
\bottomrule
\end{tabularx}}
\end{ans}
\end{keep}
\begin{dlg}
\Q{为什么 $W+R>N$ 就能读到最新值？}
\A{写成功的 $W$ 个副本和读到的 $R$ 个副本如果没有交集，就要占用 $W+R$ 个不同的副本。$W+R>N$ 时副本不够，所以两者至少有一个公共副本，它上面是最新版本，再按版本号挑出最新的值。}
\Q{Raft 和 Paxos 是什么关系？}
\A{两者都解决非拜占庭环境下的一致性问题，Raft 更容易理解。Raft 的节点分为 Leader、Follower、Candidate：Follower 收不到 Leader 的消息就变成 Candidate 发起投票，得票过半的成为 Leader，再由 Leader 向 Follower 同步日志 \kj{10-43--44}。}
\end{dlg}

% ================= 第7部分 概念题 =================
\clearpage
\section{概念题}\label{s:concept}
\begin{keep}
\begin{prob}[概念题]
理解分布式事务 ACID 四个特性，说明它们同集中式数据库的异同。此外，分布式事务的实现模型主要有哪些？
\end{prob}
\begin{ans}
原子性：事务的操作要么全做，要么全不做。分布式事务要求所有子事务要么全部提交，要么全部撤销，由 2PC 保证。
一致性：事务使数据库从一个一致状态变到另一个一致状态。分布式环境下还包括多个副本之间的一致。
隔离性：一个事务的执行不受其他事务干扰，未提交的结果不暴露给其他事务，以避免级联废弃。分布式环境下要求全局可串行，即各场地局部可串行，并且存在一致的总序。
持久性：事务提交后，修改永久保存，发生故障后可以 redo 恢复。分布式环境下还要处理场地故障和通信故障。\\
异同：分布式事务继承了集中式事务的 ACID 特性。不同在于它由多个场地上的子事务组成，需要协调各子事务的提交与撤销；除了数据的存取操作序列，还涉及大量通信原语和控制报文 \kj{7-20}。\\
实现模型：进程模型为每个子事务在相应场地建立代理者进程，另设协调者进程（根代理）负责通信与协调，只有根代理能创建子代理并决定全局提交或废弃；服务器模型在每个执行场地建立一个服务器进程，轮流为多个全局事务的子事务服务。
\end{ans}
\end{keep}

\begin{keep}
\begin{prob}[2005 年第 5.1 题]
浅析分布式数据库与集中式数据库的主要区别。
\end{prob}
\begin{ans}
数据存储：集中式数据库的数据都在一个场地；分布式数据库在逻辑上是一个整体，物理上分散在多个场地，需要做分片、分配和复制设计。
查询处理：集中式主要考虑 I/O 和 CPU 代价；分布式还要考虑网络传输代价，广域网中以传输代价为主，会用到半连接等技术。
事务管理：分布式事务分解为多个子事务，要用 2PC 保证原子性。
并发控制：分布式要保证全局可串行，还要处理分布式死锁和副本的封锁。
恢复：分布式多了场地故障和通信故障，2PC 存在阻塞问题，需要 3PC 或终结协议。
可靠性和扩展性：分布式数据库有副本，可用性和可靠性更高，容易扩展，可以并行处理；代价是系统复杂、维护开销大、安全问题更多 \kj{1-53--55}。
\end{ans}
\end{keep}

\begin{keep}
\begin{prob}[2019 年第 5 题（6 分）]
在区块链技术中，如何维护多点的数据一致性？如何理解事务特性？
\end{prob}
\begin{ans}
一致性：每个节点保存账本的完整副本，相当于全复制、对等式的分布式系统。新交易广播到全网，节点用共识算法决定由谁打包下一个区块，并让各节点接受同一个区块。公有链用 PoW、PoS 等算法，可以容忍恶意节点（拜占庭将军问题）；联盟链常用 PBFT、Raft。区块之间用哈希指针相连，篡改任何一个历史区块，后面所有区块的哈希都会对不上。出现分叉时按最长链规则收敛，所以是最终一致，交易需要等待若干个确认块才算确定。\\
事务特性：一笔交易就是一个事务。原子性：交易要么完整地打包进区块，要么不生效。一致性：交易要通过签名、余额等规则的验证，防止双花。隔离性：区块内的交易排成全序依次执行，相当于串行调度。持久性：区块被确认、后面又有足够多的区块之后，几乎无法篡改。与 2PC 相比，区块链不依赖中心协调者，用共识代替原子提交；代价是吞吐量低、确认延迟大，而且只能做到最终一致。
\end{ans}
\end{keep}

\begin{keep}
\begin{prob}[2005 年第 5.2 题、2006 年第 6 题]
通过分布式数据库的学习，你的收获是什么？\quad 分布式数据库的内容是否对你的研究方向有帮助？有何收益？说明原因。
\end{prob}
\begin{var}[作答方法]
\small 写 3 到 4 点，每点写清学到的技术，以及它在自己研究方向上的用处。可以选的技术：数据分片与分配（大规模数据按场地或特征划分存储，减少传输）；分布式查询优化（选择下推、半连接，降低通信和 I/O 代价）；分布式事务和并发控制（2PC、2PL，保证多节点上的数据正确）；复制与一致性（Paxos、Raft、NWR，在可用性和一致性之间取舍）。最后用一句话概括：数据分布在多个场地以后，带来了传输代价和一致性两类问题，课程内容基本都是围绕这两点展开的。
\end{var}
\end{keep}

\begin{var}[可能出现的其他概念]
\small
分布透明性：包括分片透明（不必知道数据属于哪个片段）、分配透明（不必知道片段在哪个场地）、局部映射透明（不必知道局部的存储形式）\kj{3-19}。\\
3PC：在 2PC 的两个阶段之间加一个“准备提交”阶段，变成投票表决、准备提交、执行三个阶段。参与者超时后可以根据其他参与者的状态决定提交或废弃，减少阻塞 \kj{7-56--64}。\\
级联废弃：一个事务读了另一个未提交事务写的数据，后者废弃时，前者也要废弃。严格 2PL 可以避免。\\
时间戳法：读操作在 $ts<wts(X)$ 时被拒绝；写操作在 $ts<rts(X)$ 或 $ts<wts(X)$ 时被拒绝；被拒绝的事务换一个新时间戳重新启动。不会发生死锁，代价是重启 \kj{9-38}。\\
CAP：一致性、可用性、分区容忍性最多同时满足两个。NWR 让用户通过参数自己选择。
\end{var}

\end{document}
